% 本文件由 LaTeX 排版 Agent 根据有效 translation.json 自动生成。
% author=Athanasius; work_id=2802; title=论圣言降生成人

\chapter{论圣言降生成人}

% structure: p0001 source=h1 target=omitted reason=page-title-already-rendered-by-parent

% structure: p0002 source=h2 target=section reason=relative-heading-rank; base=h2

\section{一、\Emph{引言。——本文的主题：圣言的谦卑与降生成人。以圣言创造说为前提。圣父藉着祂造世的同一位来拯救世界。}}

鉴于在前文中我们已——从众多要点中选取少数——充分阐述了异教徒关于偶像的错误及偶像崇拜，以及它们最初是如何被发明的；即，人出于邪恶如何为自己设计了偶像崇拜：并且藉着天主的恩宠，我们也曾稍论圣父之圣言的神性，祂的普遍眷顾与权能，以及良善的圣父藉着祂安排万物，万物因祂而运行，在祂内获得生命：那么，来吧，\Person{马卡里乌斯}（名副其实者），基督的真爱者，让我们继续我们宗教的信仰，并进一步阐明有关圣言降生成人，以及祂在我们中间的神圣显现之事——这被犹太人诋毁，希腊人嘲笑，但我们却朝拜；好叫因着圣言看似卑微的地位，你对祂的虔诚愈益增长、愈加加增。 2. 因为祂越是在不信者中被嘲笑，就越是证明祂自己的天主性；不仅因为祂亲自将人们误认为不可能之事证明为可能，而且将人们讥讽为不雅之事，以祂自己的良善装饰为雅致，又将人们凭其智慧之虚骄嘲笑为仅属人性之事，以祂自己的权能证明为神性，藉着祂被认作的屈辱——即十字架——挫败了偶像的虚假主张，又无形中赢取那些嘲笑与不信者承认祂的神性与权能。 3. 但要处理这个主题，必须回顾先前所说的；为叫你既不会不知道圣父之圣言如此崇高伟大的有形体显现的因由，也不会认为救世主穿戴肉身是祂本性的结果；而是，祂虽本性是无形的，且从起初就是圣言，却因祂自己圣父的慈爱与良善，为我们的救恩，在人身中显现给我们。 4. 因此，我们适宜以论述宇宙的创造、以及天主为其造物主来开始处理此主题，这样便可恰当地领悟：创造的更新正是那起初造化它的同一圣言的工程。因为，圣父藉着祂造世的那一位来完成救恩，似乎是毫不矛盾的。\par

% structure: p0004 source=h2 target=section reason=relative-heading-rank; base=h2

\section{二、\Emph{错误的创造观被驳斥。 (1) 伊壁鸠鲁派（偶然生成）。但身体与部分的多样性论证了一种创造理智。 (2.) 柏拉图主义者（先在质料）。但这使天主受限于人的限制，使祂不是创造者而是一个工匠。 (3) 诺斯替派（异在的巨匠造物主）。从圣经被驳斥。}}

关于宇宙的造成和万物的创造，许多人持不同观点，人人各随己意立下定论。因为有人说万物是自行产生的，出于偶然的方式；譬如伊壁鸠鲁派，他们妄自菲薄地告诉我们普遍的天主眷顾并不存在，公然违背明显的事实与经验。 2. 因为，若如他们所说，万物各依自始，与目的无关，那么必将得出万物都是单纯的存在，彼此相似而无分别。因为根据物体的统一性，必推出万物都必须是太阳或月亮，在人身上必推出全身必须是手、眼、或脚。但事实并非如此。相反，我们看到太阳、月亮和大地之间的分别；再者，在人身中，也有脚、手、头之别。这样的分离安排并非告诉我们它们是自行产生的，而是表明有一个原因先于它们；由此原因，也可能领悟天主乃是万物的制造者与安排者。\par

3. 但另一些人，包括在希腊人中享誉的柏拉图，则主张天主从在先存在的无始物质中造成了世界。因为若无材质先已存在，天主就什么也造不出；正如木匠必须先有木料在手，方能做工。\par

4. 但他们如此说，却不知自己正把软弱归给天主。因为若祂自身并非材质的因由，而只从先存的材质制造万物，祂就显为软弱，因为没有材质便不能制造任何祂所造之物；正如同木匠若无木材便不能制造任何所需之物，无疑是软弱的一样。因为，\OriginalTerm{假设地说}{ex hypothesi}，若材质本不存在，天主就什么也不会创造。那么，祂如何能被称作制造者与造物主，如果祂的制造能力依赖某种别的源头——即材质呢？因此，若果真如此，根据他们的理论，天主将仅仅是一个工匠，而不是从无中创造的创造者；也就是说，祂只加工既存的材质，却非材质自身的因由。因为，除非祂是那被造之物借以造成的材质本身的创造者，祂便在任何意义上都不能被称为创造者。\par

5. 但异端派别为自己想象了另一位不同的万物造物主，异于我们的主耶稣基督的圣父，甚至对自己所用的言语也是盲然无知。\par

6. 因为当主对犹太人说：\BibleQuote{你们没有念过：那创造者自起初就造了他们一男一女；且说，\InnerQuote{为此，人要离开父亲和母亲，依附自己的妻子，两人成为一体}？} \BibleRef{玛 19:4-5}接着，提到创造者时，说：\BibleQuote{所以，天主所结合的，人不可拆散。}\BibleRef{玛 19:6}这些人是如何断言创造与圣父无关的呢？再者，若望毫不例外地说：\BibleQuote{万物是藉着祂而造成的，且没有祂就没有什么被造成。}\BibleRef{若 1:3}那么，这位造物主怎么可能是另一位，与基督的圣父不同呢？\par

% structure: p0010 source=h2 target=section reason=relative-heading-rank; base=h2

\section{三、\Emph{真确的道理。从虚无中创造，源于天主慷慨的存在之恩赐。人被造得高于其他受造物，但无力独立坚忍。因此，天主额外超性地赐予人成为祂肖像的恩赐，并应许若在恩宠中坚忍，便得享永福。}}

他们如此徒然臆测。但敬虔的教义和按照基督的信德，却指责他们愚妄的言语为无神。因为这信德知道，宇宙并非自生，因为天主并非没有预谋；也不是由已有的物质造成，因为天主并非无能；而是从虚无中，在它毫无先前存在的情况下，天主藉着自己的圣言使之存在，正如祂首先藉着\Person{梅瑟}所说：\BibleQuote{起初，天主创造了天地。}\BibleRef{创 1:1} 其次，在极具教益的\WorkTitle{牧者}一书中说：\Quote{首先你该信天主是唯一的，祂创造并形成了万物，使它们从虚无中得以存在。} 2. \Person{保禄}也提及此事说：\BibleQuote{因着信德，我们明白诸世界是由天主的圣言构成的，这样，所见之物并不是由显然之物造成的。}\BibleRef{希 11:3} 3. 因为天主是善的，或更好说是善的本源：善者不会吝惜任何事物：因此，祂对万有不吝于存在，便藉着自己的圣言——我们的主\Person{耶稣基督}——从虚无中创造了万有。在这一切之中，祂特别怜悯了地上的人类，甚于一切，并洞悉人类由于其起源的境况，无法恒常保持稳定，于是祂赐给他们进一步的恩赐：祂并非像创造地上所有无灵动物那样，仅仅创造人，而是按照自己的肖像创造了人，并分施给他们自己圣言的某种能力；这样，他们好似具有一种圣言的反映，被造为有理性者，便能常居于真福之中，度乐园中圣人们所拥有的真正生命。 4. 但天主又知道人的意志可能摇摆于两边，便预先以一条法律和他们被安置的地点，来保障赐给他们的恩宠。因为祂将他们带入自己的园子，赐给他们一条法律：这样，如果他们遵守这恩宠，保持美善，他们就能继续在乐园中生活，没有忧苦、痛苦或挂虑，并且还有在天上获得不朽的应许；但如果他们违犯法律、转身背离，变得邪恶，他们就会知道，自己正招致本性所属的死亡中的败坏：不再生活于乐园，而是从此被逐出，死亡并常处于死亡和败坏之中。 5. 这正是\WorkTitle{圣经}所警告的，以天主的口气说：\BibleQuote{乐园中各树上的果子，你都可吃，但知善恶树上的果子，你不可吃，因为在你吃的日子，你必定死。} 但所谓\Quote{你必定死}，除了指不仅仅死去，更指常处于死亡的败坏中之外，还能有什么意思呢？\par

% structure: p0012 source=h2 target=section reason=relative-heading-rank; base=h2

\section{四、\Emph{我们的受造与天主的降生成人是最密切相关的。正如人因圣言而从虚无中被召入存在，更领受了属神生命的恩宠，同样，因一次过犯丧失了那生命，他们再次招致败坏，世界充满了无数的罪恶与痛苦。}}

你或许会感到诧异，为何我们既已立意论述圣言的降生成人，现在却讨论人类的起源。但这也正符合我们论著的目的。 2. 因为我们论及救主在我们中间的显现时，必须也论及人类的起源，好使你明白祂降来的原因正是为了我们，并明白我们的过犯唤起了圣言的慈爱，致使主既速来救助我们，又显现于人中间。 3. 因为祂降生成人，我们就是那目的；为了我们的救恩，祂如此慈爱地行事，以至显现并甚至出生于一个属人的身体内。 4. 因此，天主造了人，并愿意他常居不朽；但人却轻视并拒绝了默观天主，为自己图谋了邪恶（如前一篇论著所述），便受到他们曾被警告的死亡之惩罚；从此，他们不再保持受造时的状态，反而随自己的计谋日趋败坏；死亡便如同君王统御了他们。\BibleRef{罗 5:14} 因为违犯诫命将他们转回其本性状态，使他们就如从虚无中获得存在一样，也可预期随着时间的进程，会趋向虚无的败坏。 5. 因为，若他们是从先前的虚无常态中，因圣言的临在和慈爱而被召入存在，理所当然地，当人丧失了认识天主的知识，转回虚无（因为恶即是无，而善才是存有），他们就应——既然他们的存在来自那自有的天主——甚至永远丧失存在；换言之，他们该当解体并常处于死亡和败坏之中。 6. 因为人从本性而言是可死的，既由虚无中受造；但因着与那自有者相似（若他能通过在自己内保持这位自有者而保存这相似），他就能遏止本性败坏，保持不朽；正如\WorkTitle{智慧书}上说：\BibleQuote{遵守祂的法律，便是不死不灭的保证；}\BibleRef{智 6:18} 而既成为不朽的，他此后将如天主一般生活，我想这正与\WorkTitle{圣经}的这段话相符，圣经说：\BibleQuote{我曾说过：你们是神，你们都是至高者的儿子；但你们却要像人一样死去，像任何一位王子一样跌倒。}\par

% structure: p0014 source=h2 target=section reason=relative-heading-rank; base=h2

\section{五、因为天主不仅从虚无中造了我们；祂更藉着圣言的恩宠，白白赏赐我们一种与天主相应的生命。但人却拒绝了永恒的事物，因魔鬼的计谋转向了败坏的事物，成了自己在死亡中败坏的缘由；正如我先前所说，他们本性是可败坏的，却预定因分享圣言而来的恩宠，得以脱离本性状态，只要他们保持美善。}

2. 因为\OriginalTerm{圣言}{Word}寓居在他们中间，连他们本性的腐朽也不能临近他们，正如\WorkTitle{智慧篇}所说：\BibleQuote{其实天主造了人，原是不死不灭的，使他成为自己本性的肖像；但因魔鬼的嫉妒，死亡才进入了世界。}但当这发生后，人便开始死亡，而腐朽自此便胜过他们，甚至获得了超越其自然力量对整个种族的控制，因为它因着违犯诫命，又加上了天主的威吓作为进一步对抗他们的优势。\par

3. 因为人在恶行中，并没有停留在任何设定的界限内；而是渐渐前进，超过了所有限度：他们起初就制造了邪恶，招致了自身的死亡与腐朽；后来更转而行不义，超越一切无法无天，不只限于一种恶，而是相继发明各种新恶事，犯罪无厌。4. 因为遍地都是奸淫和偷盗，全地充满了凶杀和掠夺。至于腐朽与不义，无人顾念法律，所有罪行都到处被人施行，无论是单独还是合伙。城与城交战，国与国相攻；整个大地被内乱和战争撕裂；人人竞相在无法无天的行为上超越他人。8. 甚至违背自然的罪行也离他们不远，正如宗徒和基督的见证人所说：\BibleQuote{因为他们的女人把顺性的用处变为逆性的用处；同样，男人也是抛弃了女人的顺性用处，彼此欲火中烧，男人与男人行了丑事，就在自己身上受了这妄为当得的报应。}\par

% structure: p0017 source=h2 target=section reason=relative-heading-rank; base=h2

\section{六、\Emph{当时人类正在衰亡，天主的肖像被抹去，他的工程被毁坏。那么，要么天主必须放弃他所出的言语，那使人招致毁灭的言语；要么那曾分享\OriginalTerm{圣言}{Word}本性的，必须重新沉入毁灭，那样天主的计划就要受挫。怎么办呢？天主的良善岂能容忍此事？若是如此，当初又何必造人？那在天主方面便是无能，而非良善了。}}

为此，既然死亡战胜了人，腐朽盘踞在他们身上，人类便日渐灭亡；按天主肖像造成的理性之人正在消失，天主的工程正在瓦解。2. 因为如上所述，死亡从那时起，便依据律法上的约束控制了我们\BibleRef{创 2:15}，而这律法是不可逃避的，因为它是由天主因着违命而设定的，这事实既怪异又不合宜。3. 首先，怪异的是，天主既已发言，竟然说谎——他曾命定，人若违犯诫命，必定死去，而在违命后，人竟不死，天主的话就被破坏了。因为如果他说我们必死，人却没有死，天主就不是真实的。4. 再者，不合宜的是，受造物既曾被造为理性的，且分享了\OriginalTerm{圣言}{Word}，竟归于腐朽，经由毁灭而重返虚无。5. 因为天主所造之物由于魔鬼对人的欺骗而消亡，这不相称于天主的良善。6. 尤其是，人类中天主的杰作，或由于他们自身的疏忽，或因恶神的诡计而被除灭，是极度不适宜的。\par

7. 那么，当理性受造物正在消亡，这些工程趋于毁灭时，天主的良善该做什么呢？任凭腐朽胜过他们，让死亡紧紧抓住他们吗？那么他们起初受造又有何益处呢？因为与其受造后被遗弃于轻忽与毁灭，倒不如从未受造。8. 因为轻忽暴露了天主方面的无能，而非良善——即，他允许自己亲手所造的工程毁灭——这比从未造人更为不堪。9. 因为如果他未曾造他们，无人可指责其无能；但一旦他造了他们，从虚无中造出他们，工程却在造物主眼前毁灭，那便极其怪异了。10. 因此，绝不可将人遗弃于腐朽的洪流之中；因为这是不适宜的，也不配天主的良善。\par

% structure: p0020 source=h2 target=section reason=relative-heading-rank; base=h2

\section{七、\Emph{另一方面，还有天主本性的恒常一致，不可因我们的益处而牺牲。那么，是否应召唤人悔改呢？但悔改不能避免律法的执行；更不能救治堕落的本性。我们已招致腐朽，需要恢复到天主肖像的恩宠。唯有创造者能更新。唯有他能（1）再造万有，（2）为万有受苦，（3）代表万有到圣父面前。}}

但是，这一后果固然必然成立，但另一方面，天主的正义要求却与此相悖：即天主应当对祂就死亡所定的法则显示为信实的。因为真理之父天主，为我们的利益和保全竟成了说谎者，这实属荒谬。2. 那么，再一次，天主可能采取什么做法呢？要求人们为自己的过犯而悔改？这一点或许可说是与天主相称的；好似说，正如人们因过犯而转向了腐朽，他们藉着悔改就可以重新走上不朽之路。3. 然而悔改首先无法维护天主的正义要求。因为假如人们没有留在死亡的掌握中，祂就仍然不算更为信实；其次，悔改也不会把人们从他们的本性中召回——它仅止于阻止他们犯罪的行为。4. 可是，如果涉及的仅是一桩过失，而非随之而来的腐朽，那么悔改也就足够了。但如果，当过犯一经开始，人们便陷入了那原属他们本性的腐朽之中，并被剥夺了他们曾享有的\OriginalTerm{恩宠}{grace}，也就是那按天主的肖像而有的恩宠，那么还需要什么进一步的步骤呢？或者说，为获得这样的恩宠和这样的召回，除了那在太初从虚无中创造万有的\OriginalTerm{圣言}{Word}之外，还有什么需要呢？5. 因为，使可腐朽的成为不朽，并完整地维护圣父对于万有的正义要求，这再次正是祂的事工。因为祂身为圣父的圣言，超乎万有，唯独祂就本性而言，既能再造万有，又堪当为万有而受苦，并为万有在圣父面前作中保。\par

% structure: p0022 source=h2 target=section reason=relative-heading-rank; base=h2

\section{八、\Emph{圣言于是访问了祂虽无时不在其间的大地，并看见了这一切的邪恶。祂取了一个我们本性的身体，且是取自一位无玷的童贞女，在祂的胎中，祂使其成为自己的，为要在其中显示自己，战胜死亡，并恢复生命。}}

为此目的，那无形无体、不朽且非物质的圣言来到了我们的领域，尽管祂先前也并非远离我们\BibleRef{宗 17:27}。因为受造物中没有一部分是缺乏祂的：祂充满了各处万有，与祂的圣父同在对。但祂俯就来向我们显示慈爱，并眷顾我们。2. 祂见到有理性的受造族类正走向丧亡，死亡藉着腐朽统治着他们；又见到对过犯的威胁，给了那在我们身上的腐朽以牢固的把持，而若在法则得以应验之前就让其失效，这实属荒谬：祂再次见到，所发生之事的不妥：祂自己身为匠师所造之物正在消逝：又见到人类的极度邪恶，以及他们如何逐渐地将其加增到一个令自己难以忍受的地步：最后，见到全人类都处于死亡的刑罚之下：祂就怜悯我们的族类，体恤我们的软弱，俯就我们的腐朽，且因无法忍受死亡得胜——以免受造物丧亡，祂圣父在人们身上的手工归于虚无——便为自己取了一个身体，且是与我们同类的身体。3. 因为祂并非仅仅愿意成为有身体的，或仅仅是愿意显现出来。因为如果祂仅仅愿意显现，祂也能用其他更高妙的手段达成祂神圣的显现。然而祂取了我们同类的身体，并且不仅于此，这身体乃是由一位无玷无瑕的\OriginalTerm{童贞女}{Virgin}而来，她不认男子，一个洁净的且确实未经过人道的身体。因为祂自己是大能的，是万有的匠师，祂在童贞女胎中为自己预备了这个身体，作为自己的圣殿，使其成为祂自己的工具，在其中显示自己，并居住其中。4. 就这样，因万有都处于死亡腐朽的刑罚之下，祂从我们的身体中取了同类的一个，将其代替万有交于死亡，并奉献于圣父——祂这样做，更是出于祂的慈爱，目的是，首先，万有既被视为已在祂内死去，那导致人类毁灭的律法得以废除（因为它的能力已在主的身体上充分用尽，对于像祂同类一样的人不再有立足之地），其次，因人们曾转向腐朽，祂可藉着祂身体的归属和复活的恩宠，使他们重新转向不朽，并从死亡中复活，将死亡从他们中驱逐，如同烈火中的稻草。\par

% structure: p0024 source=h2 target=section reason=relative-heading-rank; base=h2

\section{九、\Emph{圣言，既然唯有死亡能止息这场瘟疫，便取了一个会死的身体，这身体与祂结合后，就足以代替万有，并分享祂的不朽，以阻止人类的腐朽。因祂超乎万有，祂以自己的血肉作为我们灵魂的\OriginalTerm{祭献}{offering}；因祂与我们万有合而为一，祂给我们披上了不朽。以一个比喻阐明此事。}}

因为圣言看出，人的败坏除了以死亡作为必要条件外，别无他法可以消除，而圣言作为不死不灭者，又是圣父之子，不可能遭受死亡；为此，祂取了一个能死的身体，好让它因有分于超越万有的圣言，而堪当为众人死，并且因那居住在其中的圣言而保持不朽，从此因复活的恩宠，败坏就不再能侵袭众人。因此，祂把自己所取的身体献于死亡，作为无玷的祭品和牺牲，立刻就以同等的祭献，把死亡从祂所有的同类身上除去。2. 因为天主圣言超越万有，祂自然地献出自己的殿宇和肉身的工具，作为众人的生命，以自己的死亡偿清罪债。如此，那不腐朽的天主子，因着同样的本性而与众人结合，自然地以复活的许诺，使众人穿上不朽。因为死亡中实际的腐朽，因着圣言藉着自己一个身体居住在人间，而不再能辖制人。3. 正如一位伟大的君王进入某座大城，住在其中的一所房子里，这城在任何情况下都堪受极大的尊荣，再也没有仇敌或强盗敢来侵袭、奴役它；相反，因为君王住在那里的一所房子里，它就理当受到全面的保护。万有之君王也是如此。4. 因为如今祂已来到我们的领域，住在祂同类中间的一个身体里，从此仇敌对抗人类的整个阴谋被遏止，先前占优势的死亡腐朽已被消除。若非那万有的主和救主、天主子来到我们中间，终结死亡，人类本已走向毁灭。\par

% structure: p0026 source=h2 target=section reason=relative-heading-rank; base=h2

\section{十、藉着类似的比喻，显示救赎工程的合理性。基督如何消除我们的毁灭，并以自己的教导提供了其解药。圣经中关于圣言降生成人，以及祂所完成祭献的证据。}

事实上，这项伟大工程特别适合天主的良善。1. 因为假如一位君王创建了一所房屋或一座城，若因居民疏忽而被强盗围困，他绝不会置之不理，反而要报复并收回它，视之为自己的工程，不顾居民的疏忽，只顾及自己的尊严；那么，全善圣父的天主圣言岂不更要如此，祂不会对走向败坏的人类——祂的工程——置之不理：祂既以献上自己的身体，消除了随之而来的死亡，又以自己的教导纠正了他们的疏忽，用自己的大能恢复了人所有的一切。2. 凡读到救主默感之作者著作的人，就可确信此事；他们在著作中说：\BibleQuote{因为基督的爱\BibleRef{格后 5:14}催迫着我们；因为我们如此断定：既然一个人替众人死了，那么众人就都死了；他替众人死，是为使活着的人不再为自己生活，而是为替他们死而复活了的那位生活，}我们的主耶稣基督。又说：\BibleQuote{我们却看见那位稍微逊于天使的耶稣，因受死亡之苦，被冠以尊贵和光荣；好叫他因天主的恩宠，为每个人尝到死味。}然后，他也指出为何必须是天主圣言自己降生成人的理由，如下：\BibleQuote{原来那为万物所属、万物所本的，要领许多儿子进入光荣，使救他们的元帅因受苦难而得以完全，本是合宜的；}他说这话意思是，领人类摆脱已开始的败坏，这工程非天主圣言莫属，祂起初也造了他们。而圣言之所以取了一个身体，正是为了献作如同祂自己那样的祭献，作者也提及这话：\BibleQuote{孩子既然都有血肉之躯，他也就取了同样的血肉，为能藉着死亡，摧毁那握有死亡权势的魔鬼，并释放那些因怕死而一生屈服于奴役的人。}因为藉着他自己身体的祭献，他既结束了那反对我们的律法，又藉着他赐给我们的复活希望，为我们开创了生命的新始。因为既然死亡藉着人控制了人，反过来，因天主圣言成为人，便实现了死亡的毁灭和生命的复活；正如那身负基督的人说：\BibleQuote{因为死亡既因一人而来，死者的复活也因一人而来；就如在亚当内，众人都死了，照样，在基督内，众人都要复活：}等等。如今我们不再如同服刑受死，而是如同从死者中复活的人，期待众人的普遍复活，\BibleQuote{到了预定的时期，他必显示这事的\BibleRef{弟前 6:15}}就是天主，他既成就这事，也赐给我们。这就是救主成为人的第一个原因。但也可以从下列理由看出，祂恩宠地来到我们中间，也是合宜的。\par

% structure: p0028 source=h2 target=section reason=relative-heading-rank; base=h2

\section{十一、\Emph{降生成人的第二个理由。天主知道人本性上不足以认识祂，为使人的存在有所裨益，祂赐给人关于自己的知识。祂按圣言的肖像造了他们，好让他们能认识圣言，并借着圣言认识圣父。然而，人轻蔑这事，陷入偶像崇拜，离弃了不可见的天主而投向魔术和占星术；尽管天主以多种方式启示了自己，这一切还是发生了。}}

1. 天主握有掌管万物的权能，当祂透过自己的圣言创造人类时，看到他们本性的软弱：它本身不足以认识它的创造者，甚至无法对天主形成任何概念；因为祂是非受造的，而受造物却从虚无中被造；祂是无形的，而人却以较低下的方式被塑造成肉身；而且，所造之物在各方面都远远无法领悟并认识它们的创造者——我再说，由于祂是善的，祂怜悯人类，没有让他们缺乏对自己的认识，免得他们觉得自身的存在毫无益处。\newline{}2. 受造物若不认识自己的创造者，又有何益处？或者，若不认识圣父的\OriginalTerm{圣言}{Word}（亦即理性）——他们正是在祂内领受了存活——他们如何能成为有理性者？因为如果他们只知晓地上的事，那么他们与禽兽便毫无区别。况且，若是天主不愿被他们认识，当初又何必创造他们呢？\newline{}3. 因此，为使这事不发生，祂既是善的，便使他们有分于自己的\OriginalTerm{肖像}{Image}，即我们的主耶稣基督，并按祂自己的肖像和模样造了他们：好使他们借着这恩宠，得以认识那肖像，也就是圣父的圣言，从而能经由祂而形成对圣父的概念，并在认识自己的创造者后，度幸福而真福的生活。\newline{}4. 然而，人再度因自己的乖僻，竟置这一切于不顾，藐视所赐的恩宠，如此全然弃绝了天主，使自己的灵魂昏暗，以致不仅忘记了天主的概念，还为自己发明一件又一件的虚妄。因为他们不但不按真理为自己雕刻偶像，并在生活的天主面前尊崇那些本不存在之物，\BibleQuote{敬拜事奉了受造之物，以代替造物主}\BibleRef{罗 1:25}，而且，最糟的是，他们甚至将天主的尊荣归于木石、一切物质以及人，并且变本加厉，正如我们在前一部论著中所说的。\newline{}5. 他们的不虔敬竟至如此地步，甚至去敬拜魔鬼，将其宣认为神，以满足自己的私欲。因为，如前所述，他们献上无灵的牲畜作祭品，又拿人作祭献，这正是他们当行的，就这样在那些疯狂煽动的驱使下，更紧地束缚了自己。\newline{}6. 因此，巫术也在他们中间被传授，各地的神谕误导了众人，所有人都将自己的出生与存有的影响归于星辰及一切天体，根本不曾思想可见之物以外的任何事。\newline{}7. 简而言之，事事都充满了无信仰与不法，惟有天主及其圣言无人认识，虽则祂并未在人的视线前隐藏自己，也不仅用一种方式传达对自己的认识；相反的，祂反以多种方式和形式向他们展开了这认识。\par

% structure: p0030 source=h2 target=section reason=relative-heading-rank; base=h2

\section{\Emph{12. 虽然人是受造在恩宠中，但天主预知人会遗忘，便也借着创造的化工提醒人记念祂。然而，祂还进一步设立了法律和先知，他们的职务原是为普世而设。但人却只听从自己的私欲。}}

1. 因为天主圣言的恩宠——那神圣的肖像——本身足以使人认识天主圣言，并经由祂而认识圣父；然而，天主知道人的软弱，就连他们的疏忽也预先作了安排：这样，即使他们不主动寻求认识天主，也能透过创造的化工而避免对创造者的无知。\newline{}2. 但是，由于人的疏忽逐渐沉溺于更低的事物，天主又为他们的这种软弱再次作了安排，赐下法律和先知，就是他们所熟知的人物，好叫他们即使不愿举目向天、认识自己的创造者，也能从近在身边的人获得教导。因为人从人学习高超的事，较直接。\newline{}3. 因此，他们原可以通过仰望穹苍的高处，领悟受造界的和谐，而认识其统治者，即圣父的圣言；祂以自己对万有的\OriginalTerm{天主眷顾}{providence}，使圣父为万有所认识，并为此推动万有，好使万有借着祂而认识天主。\newline{}4. 或者，如果这为他们而言太过艰难，他们至少可以遇见圣者，并透过他们学习认识天主，即万有的创造者、基督的圣父；并明白偶像崇拜实属无神，充满一切不虔敬。\newline{}5. 或者，他们甚至可以透过认识法律，而停止一切不法的行为，度修德的生活。因为法律并非只为犹太人而立，先知也并非只为犹太人而派；相反，虽然他们被派往犹太人，且遭犹太人迫害，但他们为普世而言，乃是一所认识天主、修养灵魂的圣校。\newline{}6. 尽管如此，天主的慈善、爱人之情如此浩大——然而人却被眼前的享乐及魔鬼所散播的幻象与欺诳所胜，并不肯将头抬向真理，反而给自己堆积更多的邪恶与罪恶，以致看来恍若不再有理性的样子，而从他们的行径判断，简直是毫无理性。\par

% structure: p0032 source=h2 target=section reason=relative-heading-rank; base=h2

\section{十三、那么，在这里，难道天主当保持缄默吗？容让假神窃取祂命我们向祂自己献上的敬拜吗？一位君王，若臣民反叛，在发过信函、差遣使者之后，必会亲临他们。更何况天主正要恢复我们内祂肖像的恩宠。这事，人自己是做不到的，他们不过是副本而已。因此，圣言自己必须来临（1）以再造，（2）以在肉身上毁灭死亡。}

那么，人既已如此沦为禽兽，而魔鬼的欺骗又这样遮蔽各处，隐藏了对真天主的认识，天主当做什么呢？对如此大事保持沉默，放任人被魔鬼引入歧途而不认识天主吗？2. 那么，人原先按\OriginalTerm{天主的肖像}{image of God}受造，有什么用处呢？因为对于他来说，若是受造时只成了简单的禽兽，倒比赋有理性的却活出禽兽的生活更好。3. 或者，究竟有什么必要让他一开始就领受对天主的观念呢？因为如果他连现在也不配接受，那倒不如起初就不赐给他。4. 或者，对于造了他们的天主有何利益，或于祂有何光荣，如果祂所造的人不崇拜祂，反倒以为别人是他们的创造者呢？因为这样一来，天主就显得是替别人而不是为自己造了这些人。5. 再者，仅仅一位人间的君王也不容许他所开拓的疆土归于别人，去事奉他们，或转投他人；他反而会用文书警告他们，时常借朋友派遣使者，或者必要时亲临其地，最终借亲身驾临予以谴责，只为叫他们不去事奉别人，而自己的工程不致白费。6. 那么，天主岂不更要顾惜自己的受造物，不让他们迷失而离开祂，去事奉虚无之物吗？尤其因为这样的迷失正是他们丧亡与毁灭的原因，并且既然那些曾分有天主肖像的人毁灭是不相宜的。7. 那么，天主当做什么呢？除了更新那有天主肖像的，好使人借它能再认识祂之外，还能做什么呢？可是，除非\OriginalTerm{天主真肖像}{the very Image of God}——我们的主耶稣基督——亲临，这事怎能实现呢？因为借人的方式是不可能的，既然他们只不过是按肖像受造的；借天使也不行，因为他们就连天使也不是（天主的）肖像。于是，\OriginalTerm{圣言}{Word}亲自降临了，好让祂作为圣父的肖像，能重新将那按照肖像受造的人再造出来。8. 然而，再者，除非死亡与朽坏被除去，这事也不能成就。9. 因此，祂秉自然之宜，取了一个可死的身体，好使死亡在祂身上一次永远地被除去，而使那些按祂肖像受造的人可以再度更新。因此，唯有\OriginalTerm{圣父的肖像}{Image of the Father}，为这急需已足够了。\par

% structure: p0034 source=h2 target=section reason=relative-heading-rank; base=h2

\section{十四、\Emph{一幅被污损的画像必须根据原作复原。因此，圣父的子来寻找、拯救和更新。别无他途。人既已盲目，就不能看见以得医治。受造界的见证既未能保守他，也不能使他回返。唯有圣言能作到。但如何作到？唯有借着将自己显示为人。}}

因为就如当画在木板上的肖像被外来的污渍蚀去时，那肖像的本人必须再来，好使肖像能在同一块木板上得以更新：因为，为着他的画像，就连那绘制画像的木板也不致丢弃，却在上面重新描摹出轮廓。2. 同样，圣父的至圣圣子，就是圣父的肖像，也来到我们的境地，来更新那曾按祂的肖像受造的人，并借罪过的赦免去寻找那迷失的人；如福音中祂亲口所说的：\BibleQuote{我来是为寻找并拯救迷失的人。} 因此祂也对犹太人说：\BibleQuote{人除非重生，} 意思并非如他们所想的由母胎出生，而是论及灵魂在天主肖像的模样中重生并被再造。3. 然而，既然狂热的偶像崇拜与不敬天主占据了世界，对天主的认识已被隐藏，那么教导世界认识圣父的职责属于谁呢？也许有人说是人吗？可是，人并无能力遍及天下；因为他们既无体力远行，又无法在这事上博取信任，靠自己也不足以抵御恶神的欺骗和蒙骗。4. 因为当所有人都因魔鬼的欺骗和偶像的虚幻而在灵魂上被打倒并混乱时，他们怎能争取到人的灵魂和心智——何况他们甚至看不见这些？一个人又怎能转化他看不见的东西呢？5. 但也许有人会说，受造界已经足够了；但如果受造界足够，这些大恶就根本不会发生。因为受造界已经在场，但人们照样在对天主的认识上陷于同样的错误。6. 那么，除了那洞察灵魂和心智、赐予受造万物运动、并借它们使人认识圣父的圣言之外，还需要谁呢？因为祂凭着自己的\OriginalTerm{眷顾}{Providence}和安排万物，教导人认识圣父，也是祂能更新这同一教导。7. 那么，这事如何成就呢？也许有人说，相同的途径仍如以前一样敞开，使祂可以再次借受造界的工程来显示有关圣父的真理。但这不再是可靠的方法。恰恰相反；因为人们先前已错过了看到这点，且已不再举目向上，而是向下看了。8. 因此，自然乐意有益于人，祂就以人的身份旅居于此，取一个与别人相似的身体，并借着地上的事物，即借祂身体的作为（教导他们），好使那些不愿从祂的眷顾与统御万物认识祂的人，至少能从祂真实身体所作的工程，认识那在身体内的圣言，并借着祂认识圣父。\par

% structure: p0036 source=h2 target=section reason=relative-heading-rank; base=h2

\section{十五、\Emph{这样，圣言俯就了人专注于有形之物的倾向，竟亲自取了一个身体。祂半路迎向人的一切迷信；无论人倾向于崇拜自然、人、魔鬼或亡者，祂都显示自己是这一切的主。}}

因为，如同一位关心自己门徒的仁爱老师，若某些门徒无法从高深科目中获益，便会俯就他们的水平，至少以更浅显的课程来教导他们；天主的\OriginalTerm{圣言}{Word}也是如此。正如\Person{保禄}所说：\BibleQuote{因为\BibleRef{格前 1:21}在天主的智慧中，世人凭自己的智慧，既然未能认识天主，天主就乐意以宣讲的愚拙，来拯救那些相信的人。} 2. 因为看到人拒绝了默观天主，眼睛向下，好像沉入深渊，在自然界和感官世界中寻找天主，为自己捏造出凡人和恶魔作为神明；为此，那充满仁爱、为万有的救主，天主的\OriginalTerm{圣言}{Word}，亲自取了一个身体，作为人在人中间行走，迎向所有人的感官，为要，我说，使那些以为天主是形体的人，能从主借其身体所成就的事中领悟真理，并借着祂认识父。 3. 如此，他们既为人，思想全然属人，无论将感官专注于何种对象，便在那里看见自己受到迎候，并从各处受教于真理。 4. 因为，如果他们敬畏地观看受造界，便看见受造界宣认基督是主；或者，如果他们的心思偏向人，以致认为人是神，那么从救主的工程来看，若将他们加以比较，在众人中独有救主显为天主子；因为在其他人中，未曾行过像天主的\OriginalTerm{圣言}{Word}所行的那些工程。 5. 又或者，如果他们偏向恶神，那么，看见恶神被\OriginalTerm{圣言}{Word}赶逐，他们便该知道，惟有祂，天主的\OriginalTerm{圣言}{Word}，是天主，那些灵怪并非如此。 6. 又或者，他们的心思已经堕落到死者的地步，以致崇拜英雄和诗人所言的神明，那么，看见救主的复活，他们便该明认那些是假神，惟独主是真实的，是父的\OriginalTerm{圣言}{Word}，是死亡也无法管辖的主。 7. 为此缘故，祂出生并显现为人，死亡并复活，以祂自己的工程使先前所有人的作为黯然失色，无论人的偏向在哪一方，都要从那里将人召回，教他们认识祂自己的真父，正如祂亲自说的：\BibleQuote{我来，是要拯救和寻找那迷失了的。}\par

% structure: p0038 source=h2 target=section reason=relative-heading-rank; base=h2

\section{\Emph{祂来，是要以人的身份吸引人那受感官束缚的注意力，从而引领他们认识祂是天主。}}

因为人的心灵最终已沉溺于感官之事，\OriginalTerm{圣言}{Word}便以肉身显现来隐饰自己，好能作为人，将人吸引到自己跟前，使他们的感官集中在祂身上，并且，人从此看祂为人，却因祂所行的工程而信服祂不仅是人，也是天主，是真天主的\OriginalTerm{圣言}{Word}和智慧。 2. 这也正是保禄所要指出的，他说：\BibleQuote{使你们在爱德上生根建基，能同众圣人领悟那广、长、高、深，并知道基督那超越知识之爱，为叫你们充满天主的一切圆满。} 3. 因为\OriginalTerm{圣言}{Word}到处启示自己，上天下地，深处及广处——上，在受造界；下，在\OriginalTerm{降生成人}{Incarnation}；深处，在\OriginalTerm{阴府}{Hades}；广处，在世界——万物都充满了对天主的认识。 4. 也正因如此，祂降来之后，并未立即为万有完成祂的祭献，将自己身体交付死亡并使之复活，因为那样祂就会使自己成为不可见的。但祂借着所作的事充分显示了自己，住在其中，行这些事，显这些征兆，使人知道祂不再是仅仅的人，而是天主\OriginalTerm{圣言}{Word}。 5. 因为救主借着降生成人，要完成爱的双重大工：第一，将死亡从我们中间除去，并更新我们；第二，祂虽本不可见，无形体，却借着工程彰显自己，使人知道祂是父的\OriginalTerm{圣言}{Word}，宇宙的统治者与君王。\par

% structure: p0040 source=h2 target=section reason=relative-heading-rank; base=h2

\section{\Emph{\OriginalTerm{降生成人}{Incarnation}如何不限制\OriginalTerm{圣言}{Word}的无所不在，也不减损其纯净。（太阳之喻。）}}

因为祂并非如人可能想象的，被局限在身体之内，也不是临在于身体之中时便不在其他地方；更不是祂在推动身体时，整个宇宙便失去了祂的运作和\OriginalTerm{天主眷顾}{Providence}；相反，最奇妙的是，祂既是圣言，非但不受任何事物所包容，反而是祂自己包容万有；正如祂临在于整个受造界时，祂的存有既与宇宙有别，又藉着自己的德能临在于万物之中——安排万物的秩序，并向万物和万物之中彰显祂的\OriginalTerm{天主眷顾}{Providence}，赋予每个生命和万有以生命，包容整体而不被整体所包容，唯独完完全全、在各方面存留在祂自己的圣父内——2. 因此，即使祂临在于一个人体内，并亲自赋予其生命，祂也毫无矛盾地同时赋予宇宙生命，临在于自然的一切过程中，又超乎万物整体之外；祂虽因肉身所行的工程而被认识，却也同样由宇宙的运作而显现出来。3. 灵魂的功能在于，能藉着思想活动看到自己身体之外的事物，却无法在自身身体以外运作，或以临在推动远离身体的事物。也就是说，人绝不会因想到远处的事物，就移动或改变它们；一个人即使坐在自己家中，思考天上的星体，也不会因此移动太阳或使天穹旋转。他看见它们运行并存在，却不能实际地影响它们。4. 然而，天主的圣言在祂为人的本性内并非如此；因为祂不为身体所束缚，反而是亲自支配身体，以致祂不仅在其中，更实际地临在于万物之中，同时外在于整个宇宙，唯独居住在圣父内。5. 这正是那奇妙之处：祂同时以人的方式行走，以圣言的身份赋予万有生命，又以圣子的身份与圣父同在。因此，即使童贞女生育祂时，祂也未受到任何改变；因着在肉身内，祂的光荣也未减损：相反，祂更圣化了这身体。6. 因为祂临在于宇宙之内，并不分享宇宙的本性，相反，万有都因祂而获得生命和维持。7. 连那由祂所创造、我们可见的太阳，在天空中运行，接触地上的物体时不受玷污，也不被黑暗吞噬，反而光照并清洁它们；那么，那全圣的天主圣言，太阳的造物主和主宰，更不会因在肉身中显示自己而受玷污；相反，祂既是不朽不坏的，便赋予那原本可死的身体以生命，也洁净了它：正如经上所说：\BibleQuote{他从未犯过罪，他口里也找不到诡诈。}\BibleRef{伯前 2:22}\par

% structure: p0042 source=h2 target=section reason=relative-heading-rank; base=h2

\section{十八、天主的圣言与德能如何在祂的人性行动中运作：通过驱魔、行奇迹，以及由童贞女降生。}

因此，当受默感论及此事的作者们说到祂也吃喝、被生时，你们要明白：那身体，作为身体，是被生的，并以相称其本性的食物维持；而那位与身体结合的天主、圣言自己，在安排万有的同时，也藉着在自己身体内所行的工程，显示了自己并非人，而是天主圣言。但论到祂的这些事，是因为实际吃喝、被生和受难的身体，并非属于别人，而是属于主；并且，由于祂成了人，将这些事作为人而论及祂是合宜的，为显示祂真实地、而非看起来像是地，拥有一个身体。2. 然而，正如从这些事中人们认出祂以身体临在，同样，从祂在身体内所行的工程中，祂也使人认识自己是天主子。为此，祂向不信的犹太人大声说：\BibleQuote{我若不行我父的工程，你们就不必信我；但我若行了，纵使你们不信我，也该信这些工程，好叫你们知道并明白：父在我内，我在父内。}3. 因为，正如同祂虽不可见，却透过受造界的工程而被认识；同样，祂成了人，在不可见的身体内，也可以从祂的工程中被人认识：那位能行这些事的，不是人，而是天主的\OriginalTerm{德能}{Power}和圣言。4. 因为祂命令邪魔，邪魔便被赶走，这行为不属于人，而属于天主。或者，谁看到祂治愈人类易患的疾病，还能认为祂是人而不是天主呢？祂洁净了癞病人，使瘸子行走，开启聋子的听觉，使瞎子复明，总之，从人身上驱逐了一切疾病和残弱：由这些行为，即使最普通的观察者也能看出祂的\OriginalTerm{天主性}{Godhead}。因为谁看到祂补足天生残障者所缺，使生来瞎眼的人睁开眼睛，会不明白人的本性都隶属于祂，而祂正是其\OriginalTerm{工匠}{Artificer}和\OriginalTerm{造物主}{Maker}呢？那能补足人生来所无的，显然必定也是人类自然出生的主宰。5. 因此，在起初，当祂降临我们中间时，祂从童贞女为自己形成了一个身体，这样向所有人提供了祂天主性的不小证据，因为形成这个身体的那位，同样也是其他万物的\OriginalTerm{造物主}{Maker}。谁看到一个身体仅由童贞女、未经男人而出现，会不推断出那在其中显现的，也是其他一切身体的主宰和\OriginalTerm{造物主}{Maker}呢？6. 又或者，谁看到水的实体变化转变为酒，会不明白那位行此事的，是一切水之实体的主宰和创造者呢？为此目的，祂也以海洋之主的身份行于其上，如同行走干地，为向看见的人证明祂对万有的统御。祂用少量食物喂养如此众多的人，在自己本来没有丰足之处，产生丰盛，使五饼饱饫五千人，还有十二篮剩余，岂不是将自己显示为那位眷顾万有的主吗？\par

% structure: p0044 source=h2 target=section reason=relative-heading-rank; base=h2

\section{十九、\Emph{人，生性不为所动，将要藉着那神圣的人性来受教认识天主，而整个受造界，尤其在祂的死亡中，都承认这人性所具有的天主性。}}

然而，救主如此行，原是极好的；世人既未能认识祂显露于宇宙之间的天主眷顾，亦未能辨识祂在创造中所彰示的天主性，至少可从祂肉身的作为中恢复眼目，藉着祂得以领悟对圣父的认识，正如我先前所说，由个别的事例推及祂对整体的天主眷顾。 2. 谁见过祂制伏恶魔的权柄，谁听过恶魔承认祂是他们的主，还会怀疑祂究竟是不是天主子、天主的智慧与天主的能力呢？ 3. 因为祂甚至使受造界打破了沉默：在祂死亡的那一刻，说来奇妙，更确切地说，在祂战胜死亡的真正凯旋标记——我指的是十字架——之上，整个受造界都在宣认，那位显现并由肉身受苦者，不仅是一介凡人，更是天主子、万有的救主。因为太阳隐藏了祂的面容，大地震颤，山岳崩裂：众人都惊惧不已。这些事表明，在十字架上的基督乃是天主，而一切受造都作祂的仆役，因恐惧而为其主临在作证。如此，天主圣言藉着祂的作为向人显示了自己。下一步我们必须详述祂肉身的生命终结与历程，以及祂身体死亡的性质；尤其因这是我们所信的核心，也是全人类无一例外都熟知的事：好叫你们知道，由此亦可确知基督是天主及天主子。\par

% structure: p0046 source=h2 target=section reason=relative-heading-rank; base=h2

\section{\Emph{20. 如此，除创造者外，无人能赋予不朽；除祂的肖像外，无人能恢复天主的模样；除生命外，无人能赋予生命；除圣言外，无人能教导。而祂，为偿还我们死亡的罪债，也必须为我们而死，并从坟墓中复活，作为我们的初果。因此，祂的身体必须是可死的；但祂的身体却又不可能腐朽。}}

那么，我们现已按照所能理解的程度，部分地陈述了祂以肉身显现的原因：除救主本人外，无人能将可朽者转化为不朽，因为祂在起初也从无中造了万有；除圣父的肖像外，无人能为人重新创造天主形象的模样；除我们的主耶稣基督——祂就是生命本身——外，无人能使可死的成为不死的；除那安排万有、且是唯一真实独生子的圣言外，无人能将圣父的事教导人，并摧毁偶像崇拜。 2. 然而，既然也必须清偿众人所欠的罪债：因为，如我已说过的，众人都必须死，祂正是为此特殊原因来到我们中间：为此，在藉着祂的作为证明祂的天主性之后，祂接下来为所有人献上自己的祭献，将祂的圣殿交于死亡，以代替众人，首先为使人脱离旧日的过犯、得以自由，进而显示自己比死亡更为强大，展示祂自己的身体为不朽，作为众人复活的初果。 3. 如果我们在同一主题上重复同样的话，请不要惊讶。因为我们讲论的是天主的计划，所以我们以多种形式阐明同一意义，以免看似有所遗漏，而招致处理不充分的指责：宁可受责备于重复，也好过遗漏应当记述的事项。 4. 那么，这身体既与众人有相同的本性——因为祂取得的是人的身体，虽藉前所未有的奇迹仅由童贞女所形成，然而既是可死的，也按同类的方式死亡。但因着圣言与它结合，它便不再按自己的本性屈服于腐朽，却因寓居其内的圣言而被置于腐朽的势力之外。 5. 于是，两个奇迹同时发生了：众人的死亡在主的身体内得以完成，而死亡和腐朽却因与身体结合的圣言而全然消除。因为死亡是必要的，且必须有人为众人忍受死亡，以便清偿众人所欠的罪债。 6. 因此，正如我先前所说，圣言既因是不死的而不能死，便为自己取了一个能死的身体，好作为自己的祭献代替众人，并藉着与身体的结合而为众人受苦，\BibleQuote{使那握有死亡权势的——就是魔鬼——归于无有，并解救那些因怕死而一生困于奴役的人}\BibleRef{希 2:14-15}。\par

% structure: p0048 source=h2 target=section reason=relative-heading-rank; base=h2

\section{\Emph{21. 死亡因基督的死亡而被消灭。那么，基督为何不私下死去，或以更体面的方式死去？祂并非服从自然的死亡，而必须死于他人之手。祂为何而死？不，祂正是为此目的而来，若非为此，祂便不能复活。}}

普世的救主既为我们死了，我们这些在基督内的信徒就不再像从前那样随从法律的警告而死；因为这定罪已经终止；但腐朽既已停止，并因复活的恩宠而被除去，从此我们只是按我们身体有死的本性，在天主为我们每人所定的时间分解，为能获得更好的复活。\newline{}2. 如同撒在地里的种子，我们并不因分解而灭亡，而是被种在地里，将再次复活，死亡已因救主的恩宠而被消灭。因此，真福保禄，他被立为众人复活的保证，说：\BibleQuote{这可朽坏的必须穿上不朽，这可死的必须穿上不死的；几时这可朽坏的穿上了不朽，这可死的穿上了不死的，那时就要应验经上所记载的：死亡在胜利中被吞灭了。死亡！你的刺在哪里？坟墓！你的胜利在哪里？}\newline{}3. 那么，有人或许会说，如果祂必须为众人舍弃自己的身体而交付死亡，为什么不作为人私下放下身体，而非要走到被钉死在十字架上的地步呢？因为祂以尊严放下身体，比这样忍受耻辱的死亡更为相宜。\newline{}4. 现在，请留意，我回答说，这一类的反对是否仅仅是出于人的想法，而救主所做的确实是神圣的，并且有许多理由配得上祂的天主性。首先，因为临到人身上的死亡是随从他们本性的软弱而来；由于不能持久，他们随时间而分解。因此，疾病也临到他们，他们生病并死亡。但主不是软弱的，而是天主的德能、天主的圣言和真正的生命。\newline{}5. 因此，如果祂在私下某处、在床上，按人的方式放下了身体，人们就会认为祂也是随从了本性的软弱，并且认为祂并不比其他人有更多之处。但因为祂首先就是天主的生命和圣言，其次必须完成替众人死，因此，一方面，因为祂是生命和德能，身体在祂内获得了力量；\newline{}6. 另一方面，死亡必须发生，祂不是自行取去，而是从别人手中接受；以此成全祂祭献的时机。因为让那医治他人疾病的主自己生病，本不相宜；再者，让那赐他人力量的身体自己失去力量，也不相宜。\newline{}7. 那么，祂为什么不阻止死亡，如同阻止疾病那样？正是因为祂有此身体，而不宜阻止死亡，免得复活也受到阻碍；但同时，在祂死亡之前生病也绝不相宜，免得让人以为在身体内的那位是软弱的。祂难道没有饥饿吗？是的，祂饥饿了，这是随从祂身体的属性。但祂并未因饥饿而死亡，因为有那穿戴这身体的主。因此，即使祂死了为救赎众人，祂却未见腐朽。因为身体复活时完好无损，这身体不属于别人，而正是属于生命本身。\par

% structure: p0050 source=h2 target=section reason=relative-heading-rank; base=h2

\section{二十二、\Emph{但祂为什么不从犹太人那里收回自己的身体，以保持其不死不灭呢？（1）祂不应自己死亡，也不应逃避死亡。（2）祂来是为接受他人应受的死亡，所以死亡应从外来。（3）祂的死亡必须是确定的，以确保其复活真实。另外，祂不能因病而死，以免在医治他人时被嘲笑。}}

但有人可能会说，最好避开犹太人的计谋，以完全保护自己的身体免于死亡。现在，要让这样的人知道，这对于主也是不相宜的。因为，如同天主的圣言，既是生命，不适合自己将死亡加于己身，同样，逃避由他人加给的死亡也不合适，反而应追随它直至毁灭；因此，祂自然而然地既不自己放下身体，当犹太人商议害祂时，也不逃避他们。\newline{}2. 但这并不显示圣言的软弱，相反，显示祂是救主和生命；因为祂既等待死亡以毁灭它，又急忙成就那为众人的救恩而加于祂的死亡。\newline{}3. 此外，救主来不是为完成自己的死亡，而是为完成人的死亡；因此祂没有凭自己死亡而舍弃身体\BibleRef{若 10:17}——因为祂是生命，没有死亡——而是接受了从人而来的死亡，好当这死亡临到祂自己的身体时，彻底消灭它。\newline{}4. 再者，从以下也可以看出主的身体经历此结局的合理性。主特别关心祂将要完成的肉身复活。因为祂所要做的，是显明它为战胜死亡的纪念，并向众人保证祂已实现了腐朽的消除，以及从那以后他们身体的不朽；作为这一点的保证和为所有人存留的复活的证据，祂保存了自己的身体不朽。\newline{}5. 那么，如果再一次，祂的身体生病了，圣言在众人眼前与它分离，那么，那医治他人疾病者竟让自己的工具在疾病中衰弱，便不相宜了。因为如果祂自己的圣殿在祂内生病，祂驱逐别人的疾病怎能被人相信呢\BibleRef{玛 27:42}？因为要么祂被嘲笑为不能驱逐疾病，要么若能却不为，人们也会认为祂对别人也漠不关心。\par

% structure: p0052 source=h2 target=section reason=relative-heading-rank; base=h2

\section{二十三、\Emph{公开死亡为复活教义的必要性}}

但即使祂没有任何疾病、没有任何疼痛，私下独自\BibleQuote{在一個角落}\BibleRef{宗 26:26}，或在旷野，或在房子里，或在任何地方，隐藏了自己的身体，然后突然显现，说自己已经从死者中复活了，祂在各处都会像是在讲无稽之谈\BibleRef{路 24:11}，而祂关于复活所说的话就更加不可信了，因为根本没有人为祂的死亡作证。死亡必须先于复活，因为如果没有死亡在先，就不是复活；因此，如果祂身体的死亡是在秘密中发生的，死亡既不明显，也没有发生在证人面前，祂的复活也就被隐藏，没有证据了。2. 或者，为什么当祂复活后宣扬复活时，却要让自己的死亡在秘密中发生？或者，为什么当祂在众人面前驱逐邪魔，使生来瞎眼的人复明，变水为酒，藉此使人相信祂是\OriginalTerm{天主圣言}{Word of God}，却不在众人面前彰显祂可朽的本性为不朽坏的，好使人相信祂自己就是生命？3. 或者，祂的门徒们如何能有勇气宣讲复活，如果他们不能说祂先死了？或者，当他们说先有死亡然后有复活时，如果他们没有在场作为祂死亡的证人，而那些听到他们勇敢宣讲的人就是这些证人，他们如何能被相信？因为，即便当时祂的死亡与复活是在众目睽睽之下发生的，那日的法利塞人还是不肯相信，甚至迫使那些看见复活的人否认它，那么，如果这些事是在秘密中发生的，他们定会编造多少不信的借口呢？4. 或者，死亡的终结和对死亡的胜利如何能被证明，除非在众人眼前挑战死亡，祂已经证明死亡已死，并因祂身体的不朽而永远失效？\par

% structure: p0054 source=h2 target=section reason=relative-heading-rank; base=h2

\section{二十四、\Emph{预见的进一步反对。祂没有选择自己的死亡方式；因为祂要证明自己是在所有或任何形式的死亡中的征服者：（一个好摔跤手的比喻）。那为羞辱祂而选择的死亡，反倒成了战胜死亡的战利品；此外，这死亡也保全了祂身体的完整性。}}

但我们也必须预见到其他人可能会说的话，并予以回答。因为或许有人会这样认为：如果祂的死亡必须在众人面前且有证人，以便祂复活的故事也可信，那么无论如何，祂最好为自己设计一个光荣的死亡，至少可以逃避\OriginalTerm{十字架}{Cross}的耻辱。2. 但即使这样做了，祂也会给人怀疑的根据，以为祂并非能胜过一切死亡，而只是胜过特意为祂设计的死亡；这样，对于复活，同样也会有不相信的借口。因此，死亡临到祂的身体，不是出于祂自己，而是出于敌意的计谋，好让救主无论遭受哪种死亡，祂都能彻底废除。3. 正如一位优秀的摔跤手，技艺高超且有勇气，不会自己挑选对手，免得让人怀疑他害怕某些对手，反而让旁观者选择，尤其是当他们恰巧是他的敌人时，这样，无论他们安排谁与他交手，他都能摔倒那人，并被公认为胜过所有对手；同样，万有的生命，我们的主和救主，即基督，也没有为自己的身体设计一种死亡，免得显出祂害怕某种其他死亡；祂反而在\OriginalTerm{十字架}{Cross}上接受并忍受了由他人，尤其是祂的敌人所加诸的死亡，这死亡在他们看来是可怖、可耻且不敢面对的；这样，当这死亡也被摧毁时，祂自己就被信为是生命，而死亡的能力也被彻底消灭。4. 因此，发生了一件令人惊讶和震惊的事；因为那他们认为施加为羞辱的死亡，实际上成了战胜死亡本身的胜利纪念碑。因此，祂既没有遭受\Person{若翰}那样的死亡，被斩首，也没有像\Person{依撒意亚}那样被锯成两半；为的是即使在死亡中，祂的身体仍能保持完整无缺，不给那些将要分裂教会的人留下任何借口。\par

% structure: p0056 source=h2 target=section reason=relative-heading-rank; base=h2

\section{二十五、\Emph{为什么在所有死亡中选择了十字架？(1) 祂必须为我们承受诅咒。(2) 在十字架上，祂伸开双手，将万民，犹太人和外邦人，都团聚在祂自己内。(3) 祂在属于祂自己的领域内击败了\BibleQuote{空中权势的首领}，清理了通往天堂的道路，并为我们打开了永恒之门。}}

那些教外的人为自己堆积论点的，答复到此为止。但若是我们中间也有人询问，并非出于好辩，而是出于求知的爱心，为何祂除十字架外别无他途受死，那么也当告诉他，对我们而言，除此途径并无更好，主为我们承受这事是合宜的。 2. 因为祂若亲自来承受那加在我们身上的诅咒，若非领受预定为诅咒的死亡，祂怎能\Quote{成为诅咒\BibleRef{迦 3:13}}呢？而那死亡就是十字架。因为经上正写着说：\Quote{凡挂在木头上的人是被诅咒的\BibleRef{申 21:23}。} 3. 再说，若主的死是为众人作赎价，且因祂的死，\Quote{中间隔断的墙\BibleRef{弗 2:14}}被拆毁，外邦人的呼召得以实现，祂若不钉十字架，又怎能呼召我们归向祂呢？原来只有在十字架上，人才展开双手而死。因此，主也适合担当此刑并伸开双手，好叫祂以一手拉近古时的百姓，以另一手拉近外邦人，将两者在祂自己内合而为一。 4. 这正是祂亲自说过的话，指明祂将以何种死亡赎取众人：祂说：\Quote{当我从地上被举起来时，就要吸引万人归向我\BibleRef{若 12:32}。} 5. 还有，魔鬼既是人类的仇敌，已从天上坠落，游荡在我们低层大气中，在那里管辖同他一样背命的同僚恶神，不仅借他们施行幻术，欺骗受迷惑的人，还要阻挠那些上升的人（关于此事，宗徒曾说：\BibleQuote{遵照空中权能的霸主，即现今在悖逆之子心中运行的恶神}）；但主来是为打倒魔鬼，洁净天空，并为我们开辟升天之路，正如宗徒所说：\Quote{借着幔子，就是祂的肉身\BibleRef{希 10:20}}——而这必须经由死亡——那么，除了在天空中完成的死，即十字架之外，还能以什么其他的死来实现呢？因为只有钉十字架上成全的人，是在空中死去。故此，主受此死完全恰当。 6. 这样，祂被举起，就清除了空中魔鬼和各种恶魔的邪恶，正如祂所说：\BibleQuote{我看见撒殚如闪电般从天上坠落；}并重新开启了升天之路，又如祂所说：\BibleQuote{你们这些大门，请抬高门楣；古老的门户，请高抬门扉。}其实，并不需要为圣言自己开启大门，祂本是万物之主；祂的一切化工对它们的创造者来说也并非关闭；而是我们需要，祂用祂的肉身将我们带上去。因为祂为众人将肉身献于死亡，也借此肉身再次预备了通往天上的道路。\par

% structure: p0058 source=h2 target=section reason=relative-heading-rank; base=h2

\section{二十六、\Emph{祂于第三日复活的原因。（1）不复早，否则会否认祂真实的死亡；也不（2）更迟；（a）为保全祂身体的同一性，（b）不使其门徒长久悬疑，（c）也不等到祂死亡的见证者四散，或其记忆淡忘。}}

这样，为我们而言，十字架上的死已证明是合宜且适当的，其理由在各方面都是合理的；并且可以公正地论证，众人的救恩除了十字架之外别无正当途径。因为即使如此，即使在十字架上，祂也并未让自己隐藏；相反地，当祂让受造界见证其创造者的临在时，祂却不容许自己身体的圣殿长久留在死亡中，而是只显示了它的死亡，就借着与死亡接触，在第三日将它立刻复活起来，带给祂的身体成为不可朽坏和不受苦难的，作为战胜死亡的标记与凯旋。 2. 其实，祂本可以在死亡之后立即把身体复活，显示它是活着的；但救世主由于明智的远见，没有这样做。因为若祂立即彰显复活，可能有人说祂根本没有死，或者死亡没有完全接触到祂。 3. 同样，假如祂死亡与复活之间的间隔只有两天，祂不可朽坏的光荣可能就不显明了。因此，为使身体被证实确实死了，圣言又停留了一整天，在第三日才将它不可朽坏的身体显示给众人。 4. 如此，为证明十字架上的死亡确实发生了，祂在第三日复活了自己的身体。 5. 但也免得在身体停留长时和完全腐朽之后复活，而被人怀疑，好像祂换了另一个身体——因为人可能因时间久远而怀疑所见，忘记所发生的事——为此，祂等待不超过三天；也未使那些祂曾告知复活的人长久悬疑： 6. 而是正当祂的话还在他们耳中回响，他们的眼目仍在期待，心思仍在悬疑，并且那些杀害祂的人还活在世上，尚在现场并能见证主身体的死亡时，天主子自己，在三天之后，显示出祂那曾经死去、如今却永生不朽坏的身体；并向众人显明，身体的死亡并非出自住在内的圣言的任何本性的软弱，而是为使死亡借救世主的德能在祂身上被消灭。\par

% structure: p0060 source=h2 target=section reason=relative-heading-rank; base=h2

\section{二十七、\Emph{十字架对死亡与人关系的改变。}}

因为死亡已被摧毁，十字架已成为战胜死亡之具，而死亡已无权势，确已死去，这便是一个不小的证据，更是一个明显的确据：死亡为基督的一切门徒所鄙视，且他们都向死亡进攻，不再畏惧它；反而藉着十字圣号和基督的信德，将死亡如同已死的踏在脚下。2. 古时，在救主\OriginalTerm{神圣的居留}{divine sojourn}未发生之前，死亡对圣人也是可怕的，众人为死者哀哭，如同他们灭亡了。但如今，救主既复活了祂的身体，死亡便不再可怕；因为凡相信基督的人，都把死亡践踏如同虚无，宁死也不背弃对基督的信德。因为他们确切知道，死亡时他们并非毁灭，而是实际上开始生活，并藉着复活成为不朽。3. 那曾因死亡而恶意欢跃的魔鬼，如今死亡的痛苦既已解除，它便成了唯一真正死亡者。这事的一个证据是：人未信基督之前，视死亡为可怖之物，在它面前胆怯。但一转入基督的信德和教训，他们对死亡的鄙视便如此巨大，甚至踊跃奔赴它，成为救主战胜死亡所成就的复活的见证人。他们年纪轻轻，便急于赴死，不独男人，连女人也藉着肉身的操练来对抗死亡。魔鬼已衰弱至此，甚至以前受他欺骗的妇女，如今也嘲笑他，视他如已死且瘫痪。4. 正如一个暴君被真正的君王击败，手脚被捆，所有过路的人都嘲弄他，掌击辱骂他，不再畏惧他的暴怒和凶残，因为君王已战胜他；同样，死亡既被救主在十字架上战胜并公之于众，且手脚被捆，凡在基督内的人，经过时便践踏他，并为基督作证而讥笑死亡，嘲弄他，且诵念古时所写攻击他的话：\BibleQuote{死亡啊！你的胜利在那里？阴府啊！你的刺在那里？}\BibleRef{格前 15:55}\par

% structure: p0062 source=h2 target=section reason=relative-heading-rank; base=h2

\section{二十八、\Emph{这独特的事实必须由经验检验。怀疑者应成为基督徒。}}

那么，这难道是死亡软弱的一个微小证据吗？或者，当在基督内的少男少女轻视此生并操练死亡时，这岂不是救主战胜死亡的一个微小的明证吗？2. 因为人天生惧怕死亡和身体的朽坏；但最惊人的是：凡穿戴了十字架信德的人，竟鄙视那本为可怕的自然之事，为基督的缘故而不惧死亡。3. 正如火有燃烧的本性，若有人说有一种物质不怕火烧，反而显出火的软弱——如传说中印度人的\OriginalTerm{石棉}{asbestos}那样——那么不信这说法的人，若愿加以试验，至少会在穿戴上这种防火物质并触及火之后，随即确信所赋予火的软弱：4. 或者，若有人想看看被捆的暴君，至少可以到征服者的国境和领土内，便可见那令人惧怕的人已变为软弱；照样，倘若有人即使经过这么多的证据，经过这么多在基督内成为殉道者的人，经过在基督内卓越之士天天对死亡表现的藐视之后，心中仍然怀疑死亡是否已被消除并已终结，那么，他固然理当惊奇如此大事，但切勿固执不信，也勿在面对如此显明的事实时硬心。5. 正如拥有石棉者知道火对它没有燃烧之力，又如想看被捆暴君者去到征服者的帝国，同样，不信战胜死亡的人也该接受基督的信德，归依祂的教导，就必见到死亡的软弱和战胜死亡之事。因为许多原先不信且嘲笑的人，后来都信了，并如此蔑视死亡，甚至为基督自身成为殉道者。\par

% structure: p0064 source=h2 target=section reason=relative-heading-rank; base=h2

\section{二十九、\Emph{这里便有奇妙的效果，也有一个充足的理由，即十字架，来解释它们，正如日升解释白昼。}}

现在，如果藉着十字圣号，和对基督的信仰，死亡便被践踏于脚下，那么在真理的审判台前显然就是：正是基督自己展示了战胜死亡的战利品与凯旋，并使死亡丧失了一切力量。2. 并且，如果说先前死亡是强大的并因此可怖，如今在救主的居留及他身体的死亡与复活之后，死亡便受到了轻蔑，那么显然死亡已被正是那登上十字架的基督归于无有并战胜了。3. 正如在黑夜之后太阳升起，整个大地上被它的阳光照亮，那么无论如何就不容置疑，正是太阳到处播撒了自己的光芒，驱散了黑暗并将光明赐予万物；同样，既然死亡已被轻蔑并被踏于脚下，就从救主在肉身内的救恩显现及他在十字架上的死亡发生之时起，就必然清楚显明，正是那位又在肉身内显现的救主，已将死亡归于无有，并每天都在他自己的门徒身上展示战胜死亡的标记。4. 因为当人看到那些本性软弱的人踊跃奔赴死亡，不畏惧死亡的腐朽，也不惊怕下降阴府，反而以热切的灵魂向它挑战；并且不因酷刑而退缩，反而为基督的缘故情愿舍弃地上的生命而迎向死亡，甚或当人亲眼目睹男人、妇女和幼童为着基督宗教的缘故而奔赴并跃入死亡时，谁还会这么愚昧，谁还会这么不肯相信，谁还会这么心灵残缺，竟至于看不出、推不出正是人们所见证的基督，亲自向每一个持守他的信仰并佩有十字圣号的人，供应并赐予战胜死亡的胜利，并在他们每人身上剥夺死亡的一切力量呢？5. 因为那看见蛇被踏在脚下的人，特别是知道蛇先前的凶猛，就不再怀疑蛇已死去并完全丧失了力量，除非他心灵扭曲，连身体感官都不健全。同样，谁看见狮子被孩童戏弄，还会看不见它要么死了，要么丧失了一切力量呢？6. 所以，正如用眼睛就能看见这一切的真相一样，如今死亡既被基督的信徒戏弄并轻视，就谁也不应再怀疑，也不应有人不信，死亡已被基督消灭，死亡的腐朽已被摧毁并止息。\par

% structure: p0066 source=h2 target=section reason=relative-heading-rank; base=h2

\section{三十、\Emph{复活的事实由事实证明：（1）上面所述的战胜死亡；（2）恩宠的奇迹是一位活着者的工，是天主的工；（3）如果那些神祇（如所声称的）是真实而活着的，那么摧毁他们权能的那位就更是活着的。}}

到目前为止我们所说的，便是死亡已被消灭、主的十字架是战胜它的标记这绝非小的证据。但是关于基督——万民的共同救主和真正生命——此后所实现的身体复活而成为不朽的，对于那些心灵之眼健全的人，事实的证明比推理更为清晰。2. 因为如果如我们的论证所显示的，死亡已被消灭，并且因着基督，所有的人都将死亡踏在脚下，那么祂自己首先用祂自己的身体踏碎死亡，并将它消灭，就更是如此了。如果死亡被祂杀死，那么除了祂的身体复活并被展示为战胜死亡的纪念之外，还能发生什么呢？或者，如果主的身体没有复活，死亡已被消灭这一点又怎能得以显示呢？如果有人觉得这项复活的证明不够充分，那么就让他从眼前发生的事中确信所说的话吧。3. 因为人死后便不能施展能力，他的影响力只持续到坟墓，然后就停止了；而行动以及对人的权能，只属于活着的人：就让愿意的人观察并判断，从眼见之物中承认真理吧。4. 如今既然救主在人间施行如此伟大的事，并且日复一日地看不见地劝化如此广大的人群，从各方、从居住在希腊和异域的人中，都来归向祂的信仰，并且所有的人都服从祂的教导，还会有人仍然心中疑惑，是否救主完成了复活，是否基督活着，或者祂本身就是生命呢？5. 或者是像死人一样刺痛人的良心，使他们否认祖传的律法，并向基督的教导屈膝？或者，如果祂不再活动（因为这是死人的特征），祂怎能阻止那些活动着且活着的人的行动，使行奸淫者不再行奸淫，杀人者不再杀人，作恶者不再贪婪，亵圣者从此变得虔诚？或者，如果祂没有复活而是死了，祂怎能赶逐、追逼并打倒那些不信者说他们还活着的神祇，以及他们所朝拜的邪魔呢？6. 因为哪里呼号基督的名，哪里有对祂的信仰，那里一切偶像崇拜就被废黜，一切恶灵的骗局就被揭露，任何神灵都不能忍受这名字，甚至隐约听见就逃跑并消失了。但这项工作不是死人的，而是活人的——尤其是天主的工作。7. 尤其荒谬的是，如果说被祂赶出的邪灵和被消灭的偶像还活着，而那位驱逐他们、并用自己的能力阻止它们出现，甚至被它们全都承认为天主子的祂，却是死的。\par

% structure: p0068 source=h2 target=section reason=relative-heading-rank; base=h2

\section{三十一、\Emph{如果能力是生命的标记，那么我们从偶像的无力——无论是行善还是作恶——以及基督和十字圣号的约束能力中学到什么？死亡和邪魔由此被证明已丧失了他们的主权。上述从事实而来的论证与从基督的位格而来的论证相一致。}}

但那些不信\OriginalTerm{复活}{Resurrection}的人，却给自己提供了有力的反证：因为，若按照他们的说法，基督已死，那么他们敬拜的一切邪神和神明应该能够驱除基督；但事实相反，基督证明了它们全都是死的。2. 因为，如果死人确实毫无能力，而救主却日复一日行如此众多之事：吸引人归向宗教，劝导人修德，教诲不朽，引领人对天上事物的渴望，揭示对圣父的认识，激发面对死亡的力量，向每个人显示自己，并驱除偶像崇拜的\OriginalTerm{无神}{godlessness}，而不信者的神明和邪神却一事无成，反而在基督面前显出自己的死亡，它们的荣耀沦为无能和虚幻；另一方面，藉着\OriginalTerm{十字圣号}{sign of the Cross}，一切魔法都被制止，一切巫术都被消灭，所有的偶像都被遗弃，所有放纵的享乐都被遏制，每一个人都从地上仰望天乡：那么，我们该宣称谁是死的呢？是行如此多事的基督吗？但行事不属于死者的特性。还是那毫无能力、如同无生命般躺卧的呢？这本质上正是偶像和邪神的特性，因为它们本就是死的。3. 因为天主子是\BibleRef{希 4:12}\BibleQuote{生活且行动者}，祂日复一日工作，带来万有的\OriginalTerm{救恩}{salvation}。但死亡日复一日被证明已失去一切能力，偶像和邪神被证明是死的，而不是基督，因此没有人再能怀疑祂身体的复活。4. 但是，谁若不信主身体的复活，似乎就是不明白天主的\OriginalTerm{圣言}{Word}与智慧的能力。因为，如果祂确实取了一个身体，并——如我们论证所显示的，合理一贯地——将其归为己有，主将用它做什么呢？或者说，当圣言降临其上之后，这个身体的归宿应该是什么？它不可能不死，因为它是有死的，且要为众人献于死亡：救主正是为此目的为自己塑造了它。但它不可能留在死亡中，因为它已成为生命的宫殿。因此，当它因有死而死去，又因内住的生命而复活；而复活的事实，正是由这些行为作证。\par

% structure: p0070 source=h2 target=section reason=relative-heading-rank; base=h2

\section{三十二、\Emph{可是谁曾见到祂复活，以便相信呢？不，天主常是不可见的，只能由祂的工程认识：在这里，工程正大声作证。如果你不信，看看那些信的人，从而领悟基督的\OriginalTerm{天主性}{Godhead}。魔鬼看到了这些，尽管世人眼瞎。以上论据的总结。}}

但是，如果因为祂不可见，就否认祂确实复活了，那么拒绝相信的人早该否认自然界的运行。因为天主的特性就是虽不可见，却可从祂的工程中认识，如上面已经说过的。2. 因此，如果没有工程，他们不信那未见之事，还算合理。但既然工程大声呼号并清楚显明，他们为什么故意否认那因复活而如此明显的生活呢？即使他们的理智残损，但人还能以肉身的感官看到基督不可否认的能力与\OriginalTerm{天主性}{Godhead}。3. 因为，即使是瞎子，虽看不见太阳，但若感受到太阳发出的热力，也知道太阳在高天之上。这样，我们的反对者即使至今不信，对真理仍是盲目，但至少因其他相信者而认识祂的能力，就不该否认基督的天主性和祂所完成的\OriginalTerm{复活}{Resurrection}。4. 因为显而易见的是，如果基督死了，祂就不能驱逐魔鬼、劫掠偶像；因为一个死人，邪神是不会服从的。但是，若奉祂的名就能够明显地驱逐它们，那就必证明祂没有死；尤其是，邪神能看到世人所看不见的，如果基督死了，它们就会知道，并且完全拒绝服从祂。5. 但事实上，那些不信教者所不信的，邪神却看到了——祂是\OriginalTerm{天主}{God}——因此它们逃跑并俯伏在祂脚下，就像祂在肉身时它们所喊叫的一样：\BibleQuote{我们知道祢是谁，祢是天主的圣者}；以及\BibleQuote{啊！我们与祢有什么相干，天主子？我求祢，不要磨难我。} 6. 因此，既然魔鬼承认祂，祂的工程日复一日为祂作证，就必显明——任何人都不要与真理硬碰——救主既复活了自己的身体，又是出自圣父的真\OriginalTerm{天主子}{Son of God}，是圣父的圣言、智慧和德能，在末后的时代为万有的\OriginalTerm{救恩}{salvation}而取了一个身体，教导世界认识圣父，使死亡归于虚无，并藉复活的应许将不朽赐予万有，祂首先将自己的身体复活，作为这事初果，并由\OriginalTerm{十字圣号}{sign of the Cross}的标记展示出来，作为战胜死亡及其腐朽的纪念碑。\par

% structure: p0072 source=h2 target=section reason=relative-heading-rank; base=h2

\section{三十三、\Emph{犹太人的不信和希腊人的讥嘲。前者被他们自己的圣经所驳斥。关于祂以天主及人的身份来临的预言。}}

鉴于这些事实，并且救主的身体复活与战胜死亡已清楚证实，现在让我们驳斥犹太人的不信和外邦人的嘲笑。\newline{}2. 因为在这些事上，也许正是犹太人表示怀疑，而外邦人嘲笑的地方，他们挑剔十字架的不体面，以及天主的圣言降生成人。但我们的论证将毫不犹豫地应付这两者，尤其因为我们掌握的驳斥证据昭若白昼。\newline{}3. 因为犹太人的不信可以从他们自己诵读的圣经中得到驳斥；因为这些经文，总而言之，整部受默感的圣经，针对这些事高声宣喊，正如其明确的词句充分显示的那样。因为先知们预先宣告了童贞女的奇迹和她所生的孩子，说：\BibleQuote{看，\BibleRef{玛 1:23}童贞女将怀孕生子，人将称祂的名字为厄玛奴耳，意思是：天主与我们同在。}\newline{}4. 但那位真正伟大的、而且他们相信其言语真实的梅瑟，提及救主降生成人之事，认为这预言极关重要，并确信其真实，便用以下的话记录下来：\BibleQuote{\BibleRef{户 24:5}将有一星由雅各伯升起，一杖由以色列昂扬，他必击碎摩阿布的将领。} 又说：\BibleQuote{雅各伯，你的帐幕何其华美！以色列，你的居所何其美好！如扩展的幽谷，如河边的花园，如天主所镇的帐幕，如水畔的香柏。一人必由他的苗裔生出，他必统治万民。} 还有依撒意亚：\BibleQuote{在这孩子尚未知呼父唤母之前\BibleRef{依 8:4}，大马士革的财富和撒玛黎雅的掠物，必在亚述王面前被运走。}\newline{}5. 那么，一个人将出现是在那些话里预告的。但那位要来的是万有的主，他们又如下预言说：\BibleQuote{看，\BibleRef{依 19:1}主驾着轻快的云来到埃及，埃及的偶像必在他面前战栗。} 因为从那里也有父召他回来的话，说：\BibleQuote{我从埃及召回了\BibleRef{欧 11:1}我的儿子。}\par

% structure: p0074 source=h2 target=section reason=relative-heading-rank; base=h2

\section{三十四、\Emph{关于祂受难和死亡各方面细节的预言。}}

连祂的死亡也没有被沉默掠过：相反，在神圣的圣经中提及，甚至极其明显。为使人不因缺乏对实际事件的教导而犯错，他们甚至不惮于提到祂死亡的原因——祂承受死亡不是为了自己，而是为了所有人的不死不灭与救恩，以及犹太人反对祂的计谋和他们对祂施加的侮辱。\newline{}2. 他们这样说：\BibleQuote{一个人受鞭打，深知如何忍受软弱，因他的面容转掩：他受侮辱，被轻视。他背负我们的罪过，为我们而受痛苦；我们以为他受折磨、受鞭打、受虐待；但他被刺透是因为我们的罪过，被压伤是因为我们的恶行。我们平安的惩罚临到他身上，因他的鞭伤我们得了医治。} 惊叹圣言的慈爱啊，为了我们的缘故他受羞辱，好使我们得到荣耀。经上说：\BibleQuote{我们都像羊一样迷了路，各走各自的路；但主将我们众人的罪过归于他。他受虐待，仍然谦让，总不开口，如同被牵去宰杀的羔羊，又像在剪毛人前缄默的母羊，同样他也总不开口。他受屈辱，他的审判被除去。}\newline{}3. 然后，为避免有人因他的受苦而把他设想为平常人，圣经预先防止人的猜测，并宣布为他工作的能力以及他本性与我们相比的差异，说：\BibleQuote{但谁能述说他的世代？因为他的生命从地上被夺去。他因百姓的罪过而被置于死地。我给他恶人作他的坟墓，给他富人作他的坟穴；虽然他从未行过暴虐，口里也从未放过欺诈。主愿意因他的受苦难而使他洁净。}\par

% structure: p0076 source=h2 target=section reason=relative-heading-rank; base=h2

\section{三十五、\Emph{关于十字架的预言。这些预言唯独在基督身上如何应验。}}

但或许你既听说了祂死亡的预言，还想知道关于十字架的启示。因为连这一点也没有被忽略：圣人们以极大的清晰度将其展示出来。\newline{}2. 因为首先梅瑟预言了此事，且以高声呼喊，当他说道：\BibleQuote{你们将看见你们的生命悬在你们眼前，却不相信。}\newline{}3. 接着，他之后的先知也为此作证，说：\BibleRef{耶 11:19}\BibleQuote{我像一只驯良的羔羊，被牵去宰杀，仍不知道；他们图谋恶计谋害我说：来，让我们把一棵树扔在他的饼上，将他由活人的地上铲除。}\newline{}4. 又说：\BibleQuote{他们穿透了我的双手和双脚，数尽了我的骨骼。他们分了我的外衣，为我的长衣拈阄。}\newline{}5. 被举起的死亡，且发生在木头上，这只能是十字架的死：同样，在其他任何死亡中，手脚不会被穿透，唯独在十字架上。\newline{}6. 但既然因救主居住于人间，各方各民族也开始认识天主；对于这一点，他们也没有不留记载：因为圣经中也提及了此事。因为他说：\BibleRef{依 11:10}\BibleQuote{将有叶瑟的根子，他将起来统治万民，万民都要寄望于他。}\newline{}这就是对所发生之事的一点证明。\newline{}7. 但全部圣经都充满对犹太人不信的驳斥。因为圣经中所记载的义人、圣先知和圣祖，有谁是由仅由童女所生？又有哪个女人能在没有男人的情况下孕育人类？亚伯尔不是出自亚当，哈诺客出自耶勒得，诺厄出自拉默客，亚巴郎出自特辣黑，依撒格出自亚巴郎，雅各伯出自依撒格吗？犹大不是出自雅各伯，梅瑟和亚郎出自阿默兰吗？撒慕尔不是出自厄耳卡纳，达味出自叶瑟，撒罗满出自达味，希则克雅出自阿哈次，约史雅出自阿孟，依撒意亚出自阿摩兹，耶肋米亚出自希耳克雅，厄则克耳出自步齐吗？难道他们每人不都有父亲作为其存在的根源吗？那么，谁是由仅由童女所生呢？因为先知极其重视这个征兆。\newline{}8. 又或者，谁的诞生有天空的星辰先行，向世人宣告那诞生的呢？当梅瑟诞生时，他被父母隐藏；达味甚至不被邻人所知，就连伟大的撒慕尔也不认识他，还问叶瑟还有别的儿子吗？亚巴郎也是在诞生之后才被邻人认识为大人物。但基督诞生的见证者不是人，而是来自祂所降自的高天的星辰。\par

% structure: p0078 source=h2 target=section reason=relative-heading-rank; base=h2

\section{三十六、\Emph{关于基督王权、逃往埃及等的预言}}

但哪一位曾经存在的君王，在尚无力称呼父亲或母亲之前，就执掌王权并战胜敌人呢？达味难道不是到三十岁才登上王位，撒罗满也是长大成人之后？约阿士不是七岁时才登基，而更晚的君王约史雅，年龄约七岁才接受治理？然而即使在那年龄，他们已有能力称呼父母。\newline{}2. 那么，有谁几乎在诞生之前就执掌王权并掠夺敌人呢？或者，在以色列和犹大中——让查考此事的犹太人告诉我们——曾经有过哪位这样的君王，使所有民族寄望于他并享有和平，而非与四围为敌呢？\newline{}3. 只要耶路撒冷屹立，他们之间就有无休止的战争，他们都与以色列作战；亚述人压迫他们，埃及人迫害他们，巴比伦人攻击他们；而且奇怪的是，甚至他们的邻居叙利亚人也与他们为敌。达味难道没有攻打摩阿布人，击败叙利亚人吗？约史雅防备他的邻国，希则克雅因散乃黑黎布的夸口而恐惧，阿玛肋克与梅瑟作战，阿摩黎人抵挡他，耶里哥的居民列阵对抗纳维的儿子若苏厄吗？总之，万民与以色列之间没有友谊之约。那么，这位万民将要寄望于祂的究竟是谁，值得察看。因为必须有这样一位，而先知不可能说假话。\newline{}4. 但是，有哪一位圣先知或早期圣祖，是为所有人的救恩而死在十字架上的呢？或者是谁为医治所有人而受伤并受损呢？又有哪一位义人或君王下到埃及，以致他到来时埃及的偶像就倒了呢？因为亚巴郎去过那里，但偶像崇拜同样到处盛行。梅瑟在那里诞生，但当地迷惑人的崇拜一点也不少。\par

% structure: p0080 source=h2 target=section reason=relative-heading-rank; base=h2

\section{三十七、\Emph{\BibleBook{圣咏}{} 22:16 等。祂降生与死亡的尊威。埃及的神谕与邪魔的溃败。}}

圣经所记载的人物里，有谁的手脚被刺穿，或被挂在木头上，又在十字架上为全人类的救恩而献作祭品呢？因为\Person{亚巴郎}死在床上，结束了他的一生；\Person{依撒格}和\Person{雅各伯}也是抬脚在床上死的；\Person{梅瑟}和\Person{亚郎}死在山中；\Person{达味}死在自己家中，并未遭到民众的阴谋陷害；诚然，他曾受\Person{撒乌耳}的追杀，却平安无恙。\Person{依撒意亚}被锯成两段，但未被挂在木头上。\Person{耶肋米亚}遭受羞辱，却未在定罪下死亡；\Person{厄则克耳}受苦，但不是为了民众，而是为预示民众将要遭遇的事。2. 再者，这些人纵然受苦，却仍与众人有相同的普通人性；但圣经所宣示的那位为众人受苦的，不仅被称为人，更被称为\OriginalTerm{众生之生命}{Life of all}，纵然祂在人性上与众人相似。因为经上说：\BibleQuote{你将看见}，又说：\BibleQuote{你的生命悬挂在你眼前}；又说：\BibleQuote{谁能描述祂的世代？} 原来，一切圣人的谱系都可查考，能从起初述说，知各人生自何人；然而，那身为\Emph{生命}的世代，圣经却说无可记述。3. 那么，天主的圣经说这话是指谁呢？或者，有谁能如此伟大，以致众先知预言祂如此伟大的事？如今，在圣经里找不到别人，只有那万民的共同救主，\OriginalTerm{圣言}{Word of God}，我们的主耶稣基督。因为祂是由童贞女所出，在地上显为人，祂按肉身的世代无法描述。原来，无人能说出祂肉身之父，因为祂的身体非出自男人，而仅出自童贞女；4. 所以，无人能像谱出\Person{达味}、\Person{梅瑟}及众圣祖的系谱那样，由男人追溯救主肉身的世系。因为正是祂使星辰标记祂身体的诞生；既然圣言从天降下，理当有自己的星宿也由天而来，造物之王降临时，理当被一切受造物公开承认。5. 祂生在\Place{犹太}，并有\Place{波斯}人来朝拜祂。正是祂，甚至在肉身中显现之前，就已战胜魔鬼仇敌，胜过偶像崇拜。无论如何，各地所有的外教人，都弃绝自己祖传的传统和偶像的亵渎，如今将希望寄托于基督，并登记在祂名下，这情形你们可亲眼看见。6. 因为\Place{埃及}人的亵渎从未停止过，直到万有之主，如同驾驭云彩一般，在肉身中降到那里，摧毁了偶像的迷惑，使众人归向祂自己，并藉着祂归向圣父。7. 正是祂，在太阳和一切受造物的见证下，也在那些将祂处死的人面前被钉十字架：藉着祂的死亡，救恩临到众人，一切受造物得蒙救赎。祂是\OriginalTerm{众生之生命}{Life of all}，正是祂如同羔羊，将身体献于死亡，作为替代，为救恩赐予众人，纵然犹太人不信。\par

% structure: p0082 source=h2 target=section reason=relative-heading-rank; base=h2

\section{三十八、\Emph{天主降生成人的其他明确预言。基督前所未有的奇迹。}}

如果他们以为这些证据还不够，至少该被其他理由所说服，这些理由来自他们自己所持有的神谕。因为先知们论到谁而说：\BibleQuote{我显现给那些没有寻找我的人，我被那些没有请求我的人寻见；我说：看哪，我在这里，向那未曾呼求我名的民族；我向悖逆、顶嘴的百姓伸出了我的手。} 2. 那么，有人或许会对犹太人说：那显现的是谁呢？如果是先知，就让他们说出他何时隐藏，后来又显现。再说，这是怎样的一位先知，不仅从隐匿中显现，而且还在十字架上伸出了双手？当然不是任何义人，唯有\OriginalTerm{圣言}{Word of God}，祂本性是\OriginalTerm{无形无像的}{incorporeal}，却为了我们的缘故在肉身中显现，并为众人受苦。3. 或者，如果连这一点他们也不满足，至少该被另一个证据所平息，看看它的证明力何等清晰。因为圣经上说：\BibleQuote{你们要使下垂的手强壮起来，使发抖的膝稳固；要对胆怯的人说：勇敢些，不要害怕！看哪，你们的天主要以报复来赏报你们；祂要亲自来拯救你们。那时瞎子的眼要睁开，聋子的耳朵要开启；那时瘸子要跳跃如鹿，哑巴的舌头要欢呼。} 4. 如今他们对此能说什么呢？他们究竟怎敢面对这一点呢？因为这预言不仅指出天主将在此暂居，还宣告了祂来临的征兆和时刻。因为他们将瞎子复明、瘸子行走、聋子听见、口吃者的舌头清楚说话，与将要发生的天主的降临联系在一起。那么，就让他们说说，这类征兆何时在\Place{以色列}发生过，或在\Place{犹太地区}何处有过这类事。5. \Person{纳阿曼}，一个癞病人，被洁净了，但没有聋子听见，也没有瘸子行走。\Person{厄里亚}复活了一个死人；\Person{厄里叟}也如此；但没有一个胎生瞎子复明。老实说，复活死人固然是件大事，却不如救主所行的奇事。只是，如果圣经没有略过癞病人和寡妇死儿子的例子，那么，倘若真有过瘸子行走和瞎子复明的事，记述者也不会忽略不提。所以，既然圣经对此毫无记载，显然这些事以前从未发生过。6. 那么，这些事何时发生的，岂不是在圣言亲自在肉身中来临时吗？祂何时降临，岂不是在瘸子行走、口吃者说话清楚、聋子听见、胎生瞎子复明之时吗？这正是当时亲眼目睹的犹太人所说的话，因为他们从未听说过这些事在其他任何时候发生过：\BibleQuote{自创世以来，从未听说过有人开了胎生瞎子的眼睛。这人若不是从天主来的，什么也不能做。}\par

% structure: p0084 source=h2 target=section reason=relative-heading-rank; base=h2

\section{三十九、\Emph{难道你们还在期待另一位吗？但达尼尔预言了确切的时间。对此异议的驳斥。}}

但也许，即使他们也无法不断对抗明显的事实，他们虽然不否认所记载的，却坚称他们仍在期待这些事，并说\OriginalTerm{天主圣言}{Word of God}尚未到来。他们反复纠缠这一点，面对明证也不以厚颜为耻。 2. 但在这点上，他们将特别被驳斥，不是由我们，而是由最明智的\Person{达尼尔}，他既指出了确切的时间，又指出了救世主的神圣居留，说：\BibleQuote{七十个星期已为你的人民和你的圣城裁定了，为终止过犯，封住罪恶，除尽不义，赎清罪孽，带来永久的正义，应验神视和预言，并傅油万圣之圣；你当知道并明白，从发出命令重建\Place{耶路撒冷}，直到受傅的君王。} 3. 也许对于其他的（预言），他们或许能找到借口，将所记载的推迟到将来。但对这个他们能说什么，或能如何面对呢？这里不仅提及基督，而且声明那受傅者不单是人，而是万圣之圣；\Place{耶路撒冷}要存留直到他来，此后，在\Place{以色列}中，先知和神视都要停止。 4. 古时\Person{达味}受傅，\Person{撒罗满}和\Person{希则克雅}也受傅；但那时，\Place{耶路撒冷}和圣地仍存留，先知仍在说预言：\Person{加得}、\Person{阿撒夫}和\Person{纳堂}；后来还有\Person{依撒意亚}、\Person{欧瑟亚}、\Person{亚毛斯}及其他。再者，那些实际受傅的人被称为圣者，而不是万圣之圣。 5. 但如果他们以被掳为挡箭牌，说因此\Place{耶路撒冷}不再存在，那么关于先知他们又能说什么呢？事实上当百姓初次下到\Place{巴比伦}时，\Person{达尼尔}和\Person{耶肋米亚}在那里，\Person{厄则克耳}、\Person{哈盖}和\Person{匝加利亚}都在说预言。\par

% structure: p0086 source=h2 target=section reason=relative-heading-rank; base=h2

\section{四十、\Emph{论证（1）从预言的消失和耶路撒冷的毁灭，（2）从外邦人的归化，而且是归向梅瑟的天主。默西亚还有什么要做的，是基督没有做的呢？}}

所以犹太人在妄想，而他们所指向将来的那段时间实已来到。因为先知和神视何时离开以色列，难道不是在基督——万圣之圣——来临之时吗？因为这是一个标记，且是\OriginalTerm{天主圣言}{Word of God}来临的重要证据，即\Place{耶路撒冷}不再存在，也没有任何先知兴起或神视启示给他们——这实在是非常自然的。 2. 因为他所预表的那位已经来到，哪里还需要什么来预示他呢？真理既已临在，哪里还需要影像呢？这原是先知所以预言的原因——即直到真正的正义来临，以及那位要赎清众人罪过者到来。这就是为什么\Place{耶路撒冷}得以存留直到那时——即让他们得以在预像中练习，为那实体作准备。 3. 所以当万圣之圣来到时，很自然地神视和预言就被封存，\Place{耶路撒冷}的王国也终结了。因为君王在他们中间受傅，只持续到万圣之圣受傅为止；\Person{雅各伯}也预言说犹太人的国权将持续到他来临，如下所说：\BibleQuote{权杖不离\Person{犹大}，统帅不离他股间，直到那应得权杖者来到；万民都要归顺他。}\BibleRef{创 49:10} 4. 因此救主自己也大声宣告说：\BibleQuote{法律和先知讲说预言，直到\Person{若翰}。} 如果现在犹太人当中仍有君王或先知或神视，那么他们否认已经来临的基督就有道理了。但如果没有君王也没有神视，反而从那时起所有的预言都被封存，城和圣殿都被夺去，他们为什么如此不敬和刚愎，看到已经发生的事，却仍否认造成这一切的基督呢？又为什么当他们看到连外邦人也离弃了他们的偶像，通过基督将希望寄托于\Place{以色列}的天主时，他们却否认那位按肉身出自\Person{叶瑟}根苗，从此为君王的基督呢？因为如果万民朝拜的是另外一位神，而不宣认\Person{亚巴郎}、\Person{依撒格}、\Person{雅各伯}和\Person{梅瑟}的天主，那么他们声称天主没有来临还有道理。 5. 但如果外邦人崇敬的是那位赐给\Person{梅瑟}法律并向\Person{亚巴郎}许下承诺的同一天主，而他的言语犹太人不尊敬——他们为什么不知道，或更好说为什么故意忽略，圣经所预言的主已经光照世界，并以可见的形体显现于世，正如圣经所说：\BibleQuote{主天主已光照我们；} 又说：\BibleQuote{他发出圣言，医治了他们；} 又说：\BibleQuote{不是使者，不是天使，而是主亲自拯救了他们？} 6. 他们的状况好比一个精神失常的人，看到太阳照亮大地，却否认那照亮大地的太阳。因为那一位他们期待当他来时，他还要做什么呢？召叫外邦人吗？他们已被召叫了。使预言、君王和神视停止吗？这已经发生了。揭露偶像崇拜的不虔吗？这已被揭露并被定罪了。或是毁灭死亡？死亡已经被毁灭了。 7. 那么还有什么没有发生，是基督必须做的呢？还有什么未应验，使得犹太人如今可以不信而无罪呢？因为如果，我说——正如我们实际看到的——在他们当中再也没有君王、先知、\Place{耶路撒冷}、祭献和神视，而整个大地却充满对天主的认识，外邦人离弃他们的不虔，如今借着圣言，即我们的主耶稣基督，投奔到\Person{亚巴郎}的天主面前，那么即使对最顽固的人，也必然是显然的：基督已经来临，他用自己的光明普照众人，并将关于他圣父的真道和属神的教导赐给了他们。 8. 这样，人们可以合情合理地由这些和其它来自圣经的论证来驳斥犹太人。\par

% structure: p0088 source=h2 target=section reason=relative-heading-rank; base=h2

\section{四十一、\Emph{回答希腊人。他们承认\OriginalTerm{圣言}{Logos}吗？如果祂在宇宙的有机体中显现自己，为何不在一具身体中呢？因为人体是同一整体的部分。}}

但人们不得不对外邦人感到极为惊讶，他们嘲笑不当嘲笑之事，却对自己立起的木石形状的羞耻麻木不仁，看不到自己的耻辱。2. 不过，既然我们的论证不乏证明，来，让我们也以合理的理由使他们羞愧——主要根据我们自己也看到的事。因为我们这边有什么可笑或值得嘲弄的呢？难道仅仅因为说\OriginalTerm{圣言}{Word}在身体中显现吗？但若他们显为真理之友，甚至他们也会承认这绝非荒唐。3. 如果他们根本否认有天主圣言的存在，那就是肆意嘲笑他们所不知的。4. 但若他们承认有天主圣言，祂是宇宙的统治者，圣父在祂内造成了造化，万物因祂的天主眷顾获得光、生命和存在，祂统御万有，由此从祂天主眷顾的化工中得以认识祂，并藉着祂认识圣父——请思量，他们岂不是无意中自嘲。5. 希腊的哲学家说宇宙是一巨大的身体；诚然如此。因为我们看见它及其各部分作为感官的对象。那么，如果圣言在宇宙这身体中，并已结合于整体及所有部分，我们说祂也结合于人，有何出奇或荒唐呢？6. 因为如果祂在任何身体中都属荒唐，那么祂与整体结合，并以祂的天主眷顾赋予万物光与运动，也属荒唐。因为整体亦是身体。7. 但若祂与宇宙结合，并在整体中被认识是合宜的，那么祂亦在人体中出现，并藉祂而受光照与运作，也必定是合宜的。因为人类像其余部分一样，是整体的一部分。而如果一个部分被采用作为祂的工具来教导人认识祂的天主性是不合宜的，那么祂借整个宇宙来彰显自己，也必最为荒唐。\par

% structure: p0090 source=h2 target=section reason=relative-heading-rank; base=h2

\section{四十二、\Emph{祂与身体的结合基于祂与整个造化之关系。祂使用了一具人体，因为祂正是要以此向人启示自己。}}

因为正如，整个身体由人赋予生命与光照，假如有人说，人的能力不应也存在于脚趾，那会被视为愚昧；因为虽承认人遍在于整个身体且在其中运作，却拒绝其亦在部分中；同样，凡承认并相信\OriginalTerm{圣言}{Word}在整个宇宙中，且整体由祂光照与推动者，也不应认为一具人体由祂获得运动与光照是荒唐的。2. 但如果他们因为人类是受造的，且从虚无中被造，就认为我们所说的救主在人中的显现是不合宜的，那么他们也早该把祂从造化中驱逐出去；因为造化也是由圣言从虚无中产生的。3. 但若纵然受造物是制成的，圣言在其中并不荒唐，那么祂在人内亦不荒唐。因为他们对整体形成任何观念，必然也要应用于部分。因为如前所述，人也是整体的一部分。4. 因此，圣言在人内毫无不合宜之处，同时万物都由祂获得光、运动和生命，正如他们的作者也说：\BibleQuote{我们生活、行动、存在，都在他内}，\BibleRef{宗 17:28}。5. 那么，我们所说的有何可嘲笑的呢？如果圣言使用祂所在之物作为工具来彰显自己？因为如果祂不在其中，祂也不能使用它；但我们既然已承认祂在整体及部分中，那么祂在所在之处彰显自己，又有何不可信呢？6. 因为凭祂自己的能力，祂完全结合于每一个与全部，毫无保留地安排万物，所以若祂愿意，无论藉太阳、月亮、天穹、大地、水或火，来说话并显示自己和圣父，无人能称为不当；因为祂同时将万物容纳为一，且事实上不但在万有中，也在所述的部分中，并在那里无形地彰显自己。同样，若祂安排整体，赋予万物生命，且愿意通过人来彰显自己，祂使用一具人体作为工具以彰显圣父的真理与知识，这也绝非荒唐。因为人类本身也是整体的一部分。7. 如同\OriginalTerm{理智}{Mind}遍及整个人，却由身体的一部分，我指舌头，来表达，我想没有人会说理智的本质因此被贬低；同样，若遍及万物的圣言使用了人的工具，这也无法显得不合宜。因为，如我之前所说，若使用身体作为工具是不合宜的，那么祂在整体中也就不合宜了。\par

% structure: p0092 source=h2 target=section reason=relative-heading-rank; base=h2

\section{四十三、\Emph{祂以人的形态而非更尊贵的形式来临，因为（一）祂来是为拯救，而非为炫耀；（二）在受造物中，唯独人犯了罪。由于人不肯借宇宙中的作为认识祂，于是祂以人的身份来到他们当中，在他们自我局限的范围内施行作为。}}

如果有人问：那么，祂为何不通过受造界中其他更为尊贵的部分显现，并使用某种更为尊贵的工具，比如太阳、月亮、星辰、烈火或空气，而单单只用人呢？须知，主来不是为了炫耀，而是为医治与教导那些受苦的人。2. 因为，想要炫耀者的方式，无非是显现自己，使观看者目眩；但对于寻求医治与教导者的方式，则不单是寄居于此，而是献身援助需要者，并以需要祂的人所能承受的方式显现；以免因超越受苦者的承受能力，反而困扰那些需要祂的人，使天主的显现对他们毫无裨益。3. 在受造物中，除了人以外，没有谁对天主的观念陷入歧途。原来，太阳、月亮、穹苍、星辰、水、空气，无一偏离其秩序；它们认识自己的\OriginalTerm{造化者}{Artificer}与\OriginalTerm{主宰}{Sovereign}——\OriginalTerm{圣言}{Word}，便仍保持受造时的本相。唯独人，弃绝了善，反以虚无之物代替真理，将当归于天主的尊崇以及对祂的认识，归给了魔鬼和具有石头形状的人。4. 因此，这是合理的：既然忽略如此重大的事，与\OriginalTerm{天主的美善}{Divine Goodness}不相称，而人又无法在整体中认出祂是万有的安排者与引导者，祂便取了整体的一部分，即祂的人性身体，作为工具，并与之结合；使那些无法在整体中认出祂的人，至少能在部分中认识祂；使那些无法仰望祂无形大能的人，至少能从相似于自己的形体出发，推想并观瞻祂。5. 正因为他们是人，才能藉着与其本性相同的身体，以及藉这身体所行的神圣工程，更快更直接地认识祂的父；通过比较判断，这些工程并非出自人，而是天主藉祂而行的工程。6. 若如他们所言，圣言藉身体的工作得以被认识是荒谬的，那么祂藉宇宙的工程得以被认识，也同样荒谬。因为，正如祂在受造界内，却丝毫未分享受造界的本性，反倒是万有分享祂的大能；同样，当祂以身体为工具时，并未分享任何肉身的属性，相反，祂自己甚至圣化了那身体。7. 因为，甚至备受希腊人尊崇的\Person{柏拉图}也说，宇宙的作者看到宇宙在暴风雨中颠簸，有坠入混沌深渊的危险，便坐到灵魂的舵轮旁，前来救援，纠正它的一切灾祸；那么，我们所说的：人类陷入错谬时，圣言降临其上，显现为人，为要在风暴中藉其引导和美善拯救人类，又有什么不可信的呢？\par

% structure: p0094 source=h2 target=section reason=relative-heading-rank; base=h2

\section{\Emph{既然天主以一言造人，为何不也以一言恢复人？然而（1）从无中创造不同于修复已存在之物。（2）人因有特定的需求，就需要特定的救治。死亡已深植于人性中：因此祂必须将生命紧密地与人性交织在一起。所以\OriginalTerm{圣言}{Word}\OriginalTerm{降生成人}{Incarnate}，为要在那被篡夺的领域迎战并战胜死亡。（稻草与石棉的比喻）}}

然而，也许他们因惭愧而同意这一点，却会说：天主若愿意改革并拯救人类，只须发出一道命令，无需祂的圣言取肉身，正如祂从前从虚无中创造他们时所作的那样。2. 对他们的这种异议，一个合理的答复是：从前，当万物尚不存在时，创造一切所需要的，只是一道命令和纯粹的意志。然而，当人已被造成，需要治疗时——不是对不存在之物，而是对已形成之物——自然的结果就是，这医师和救主应在已形成之物中显现，以便治疗那已存在之物。为此，祂成了人，并以祂的肉身作为人性工具。3. 因为如果这不是正确的方法，那么选择使用一个工具的圣言，又该如何显现呢？祂又该从哪里取这工具，难道不是从那些已经存在、且因自己同类而需要祂的天主性者中取吗？因为需要救恩的，并非不存在之物，只需一道命令即可，而是已经存在的人，正走向腐朽与败坏。因此，圣言使用人性工具并在各处启示自己，乃是自然且正当的。4. 其次，你们也必须知道：那临到的腐朽，并非在肉身之外，而是与肉身紧密相连；因此，必须以生命取代腐朽，依附于肉身；这样，正如死亡在肉身内产生，生命也在其中产生。5. 如果死亡在肉身之外，那么生命同样在肉身之外产生才是合适的。但既然死亡紧密缠绕肉身，如同与它合为一体般辖制它，那么生命也必须紧密缠绕肉身，这样，肉身因穿戴生命而取代死亡，得以摆脱腐朽。此外，即使假设圣言来到肉身之外，而非在其内，死亡固然会被祂战胜，完全合乎本性，因为死亡对生命毫无权能；但依附于肉身的腐朽仍会存留在内。6. 因此，救主合理地穿上了肉身，为的是肉身与生命紧密相连后，不再作为可朽之物滞留于死亡中，而是穿戴了不死不灭，从此复活并常存不朽。因为肉身既已穿戴腐朽，除非穿戴生命，否则无法复活。同样，死亡按其本性，也不能在肉身之外显现。因此，祂穿上了肉身，好使死亡在肉身中被找到，并被抹除。若非使有死之物获得生命，主又怎能被证明就是那生命呢？7. 正如枯草本易被火烧毁。假设有人首先使火远离枯草，虽然枯草未被烧着，但它仍旧不过是枯草，仍畏惧火的威胁——因为火的本性就是焚毁它；但若有人其次，用大量石绵包裹它——石绵据说是抵抗火的物质——枯草便不再惧怕火，因被不可燃物质包裹而安全；8. 同样，就肉身与死亡而言，我们可以说，如果死亡只因祂的一道命令便远离肉身，肉身仍按其本性是可朽与可腐朽的；但为使这情形不发生，肉身穿戴了\OriginalTerm{无形无像}{incorporeal}的天主圣言，因而不再畏惧死亡或腐朽，因为它以生命为衣，腐朽在其中已被消灭。\par

% structure: p0096 source=h2 target=section reason=relative-heading-rank; base=h2

\section{四十五、\Emph{这样，受造界的每一部分再次彰显了天主的荣耀。自然作为其创造者的见证者，借着奇迹，为降生成人的天主献上了第二份证词。自然的见证虽被人类的罪所扭曲，却因此被强行带回了真理。若这些理由还不够，就让希腊人看看事实吧。}}

因此，天主圣言前后一致地取了肉身，并使用人性工具，为的是也赋予肉身生命，并且如同在创造中借其工程被认识一样，也在人内工作，在各处彰显自己，不留下任何不含祂的天主性和对祂的认识的空白。2. 我要从前面所说的重新说起，救主这样做，是为了如同祂以临在充满万有一切方所，也同样能以对自己的认识充满万有，正如圣经所说：\BibleQuote{大地充满了主的认识。}3. 因为，若有人仰望高天，便可见其秩序；若他不能仰面向天，只能看向人，他也能看见祂的能力——远超凡人能力——通过祂的工程彰显，从而得知在众人中唯独祂是天主圣言。或者，若有人误入魔鬼之中，惧害它们，他可见此人将它们驱除，从而认定祂是它们的主宰。又或者，若有人沉溺于水，以为水即是神——如埃及人崇拜水一般——他可见水因祂而改变性质，从而知晓主是众水的创造者。4. 若有人甚至下到阴府，敬畏那些曾下到那里的英雄，视他们为神，他却可见基督复活和战胜死亡的事实，从而推知，在他们中间，唯有基督是真天主和主。5. 因为主触摸了受造界的每一角落，将一切从所有幻象中释放和唤醒；正如保禄所说：\BibleQuote{解除了率领者和掌权者，在十字架上凯旋而归：}\BibleRef{哥 2:15}为使任何人都不再可能受骗，而是到处都能找到真正的天主圣言。6. 这样，人四面被围，注视圣言的天主性处处彰显——即，在天上、在阴府、在人内、在地上——便不再易受关于天主的欺骗，反而当唯独敬拜基督，并借祂正确地认识圣父。7. 因此，通过这些基于理性的论证，外邦人反过来应被我们合理地说服。但是，如果他们觉得这些论证不足以使他们惭愧，至少让他们因那些众人皆可见的事实，确信我们所说的这一切。\par

% structure: p0098 source=h2 target=section reason=relative-heading-rank; base=h2

\section{\Emph{自\OriginalTerm{降生成人}{Incarnation}起，偶像崇拜、神谕、神话、邪魔力量、巫术和外邦哲学的失信。旧宗教崇拜严格局限于本地且独立，而基督的崇拜则是至公且一致的。}}

当人们何时开始放弃偶像崇拜，岂不是自天主、天主的真\OriginalTerm{圣言}{Word}来到人间之时？希腊各地以及各处的神谕又是何时废弛并归于虚空，岂不是自救主显现于大地之时？2. 当那些在诗人笔下被称为神祇和英雄的开始被证实不过是必死的人，岂不是自主建立起对死亡的战胜，并持守他所取的肉体不朽，从死者中复活之时？3. 当邪魔的欺诈和疯狂何时开始遭人鄙弃，岂不是当天主的能力——圣言，也是这一切的主宰——因人的软弱而屈尊，显现于大地之时？当巫术的技巧和学派何时开始被践踏，岂不是当圣言的天主性显现于人间之时？4. 简言之，希腊人的智慧何时变为愚拙，岂不是当天主的真智慧显于大地之时？因为从前整个世界和每个地方都被偶像崇拜引入歧途，人们别无他神，只以偶像为神。但现在，全世界的人们都在离弃偶像的迷信，投奔基督；且因敬拜他为天主，就藉着他得以认识他们素不认识的父。5. 而奇妙的是：从前崇拜的对象种类繁多、数目庞大，每个地方各有自己的偶像，而那在他们中被视为神的，连越过边界到邻近之处、说服邻近的人民敬拜他也无能为力，甚至在自己的人民中也几乎无人服事；因为没有人敬拜邻人的神——相反地，各人持守自己的偶像，以为它是万有的主宰——唯独基督在所有民族中被敬拜为同一的主；偶像的软弱所不能的——即使说服那些住在邻近的人——基督却成就了，他不仅说服了邻近的人，而且说服了全世界，要敬拜同一的主，并藉着他敬拜天主，即他的父。\par

% structure: p0100 source=h2 target=section reason=relative-heading-rank; base=h2

\section{\Emph{众多神谕——在圣地想像的显现等——均被十字圣号驱散。旧神明被证实不过是人。巫术暴露。哲学只能说服少数精选的本地小团体关于不朽与善——而悟性不高的人却将\OriginalTerm{超性生命}{supernatural life}的原理注入教会的群众中。}}

从前各处充满了神谕的欺骗，就如在\Place{德尔斐}、\Place{多多纳}、以及\Place{彼俄底亚}、\Place{吕基亚}、\Place{利比亚}、\Place{埃及}，还有\Place{卡比里}的神谕和皮同女祭司，都因人的幻想而被人看重；但现在，自从基督开始在各处被宣讲，他们的疯狂也停止了，他们中间再没有能预言的了。2. 从前邪魔欺骗人的幻想，附在泉水、河流、树木或石头上，用戏法迷惑无知者；现在，因圣言的神圣临在，它们的欺骗已止息。因为人只要使用\OriginalTerm{十字圣号}{Sign of the Cross}，就能驱除它们的欺骗。3. 从前人把诗人所提及的\Person{宙斯}、\Person{克罗诺斯}、\Person{阿波罗}以及诸英雄当作神明，妄加崇拜；现在救主已显现于人间，那些神明都被揭露为必死的人，唯独基督在人间被认作真天主，天主的圣言。4. 至于他们中间所推重的巫术，又怎么说呢？在圣言寄居我们中间之前，它在埃及人、迦勒底人和印度人中是很强大而活跃的，令看见的人惊惧；但由于真理的临在，以及圣言的显现，它也被彻底驳倒，完全归于无有。5. 至于外邦智慧和哲学家的虚夸论点，我想无需我们举证，因为这奇迹就摆在众人眼前：希腊的博士们虽写了那么多，却不能说服他们本地的少数人相信不朽和德行的生活；唯独基督，用普通的言词，并藉着不擅长言辞的人，竟在全世界说服了充满众人的整个教会，使人轻看死亡，思念不朽的事；轻看暂时的事，转眼注视永恒的事；不顾地上的荣耀，只追求天上之事。\par

% structure: p0102 source=h2 target=section reason=relative-heading-rank; base=h2

\section{\Emph{更多事实。基督徒童贞者和苦修者的节德。殉道者。十字圣号对抗邪魔和巫术的能力。基督藉其大能显示自己超过一个人，超过一个巫师，超过一个神灵。因为这一切都完全属他管辖。因此他是天主的圣言。}}

我们这些论证并非仅是空谈，而是有实际经验作为其真理的见证。\newline{}2. 谁若愿意，就上前去，在基督的贞女们和实践圣洁贞洁的青年身上，看看德行的明证；在祂如此众多的殉道者队伍中，看看永生的保证。\newline{}3. 谁想凭经验验证我们刚说的话，就让他在魔鬼的欺骗、神谕的诈欺和魔术的奇事面前，使用那在他们中间被嗤笑的十字架的标记，他必看见，藉此标记，魔鬼逃遁，神谕止息，一切魔术巫术归于无有。\newline{}4. 那么，这位基督是谁，何等伟大？祂凭自己的名号和临在，使一切事物在各方面黯然失色、归于无有，唯独祂对抗一切时强大无比，并以自己的教导充满了整个世界。让那些喜欢嗤笑而不知羞耻的希腊人来告诉我们吧。\newline{}5. 如果祂只是人，那么一个人怎能超越他们自己敢称为神的那些人的能力，并用祂自己的能力证明他们一无所有呢？但如果他们称祂为术士，那么一个术士又怎能毁灭一切魔术，而不是使之巩固呢？因为如果祂只战胜了个别的术士，或仅压倒其中一个，那么他们便可以说祂是以更优越的技巧胜过其余；\newline{}6. 但如果祂的十字架确实完全战胜了一切魔术，甚至战胜了魔术之名，那么显然救主不是术士，因为连其他术士所召请的魔鬼也都逃避祂如同逃避自己的主一样。\newline{}7. 那么祂是谁，让那些专以玩笑为事的希腊人来告诉我们吧。或许他们会说祂也是一个魔鬼，因此有力量。但他们尽管这样说，嘲笑必归于他们自己，因为他们可再次因我们先前的证据而蒙羞。因为一个驱除魔鬼的人，怎可能是魔鬼呢？\newline{}8. 如果祂只是驱逐个别的魔鬼，那么他们也许可以说祂是凭魔王战胜了小鬼，正如犹太人想侮辱祂时对祂所说的那样。但如果一提起祂的名字，一切魔鬼的疯狂便被根除驱逐，那么显然他们在这里也错了，我们的主和救主基督并非如他们所想的某种魔鬼的能力。\newline{}9. 这样，如果救主既不是纯粹的人，也不是术士，也不是什么魔鬼，而是凭祂自己的\OriginalTerm{天主性}{Godhead}，将使诗人们所传的教义、魔鬼的幻妄和外邦人的智慧尽都归于无有、使之黯然失色，那么显而易见，人人必将承认，祂就是真天主子，就是从起初就有的\OriginalTerm{圣父}{Father}的\OriginalTerm{圣言}{Word}、智慧与能力。因为这也是为什么祂的工程绝非人的工程，而是公认超乎人之上，确实是天主的工程，这无论是从事实本身，还是从与〔其余〕人类的比较中，都可以看出来。\par

% structure: p0104 source=h2 target=section reason=relative-heading-rank; base=h2

\section{\Emph{49. 祂的诞生与奇迹。你们因\Person{阿斯克勒庇俄斯}、\Person{赫拉克勒斯}和\Person{狄俄尼索斯}的作为而称他们为神。将他们的作为与祂的对比，以及祂死亡时的奇事，等等。}}

因为何曾有人，仅从一位贞女，为自己形成了身体呢？又何曾有人，像万有的共同之主那样，医治过这样的疾病呢？或者谁弥补了人性的缺失，使一个生来瞎眼的人得以看见呢？\newline{}2. \Person{阿斯克勒庇俄斯}在他们中间被尊为神，因为他行医，并为病体寻找药草；他并非自己从土中造出它们，而是凭从自然中得出的学问发现它们。但这比起救主所做的，又算什么呢？救主不是治好一处创伤，而是改变了人原来的本性，恢复了健全的身体。\newline{}3. \Person{赫拉克勒斯}在希腊人中被当作神来敬拜，因为他与自己的同类人作战，并用计谋消灭了野兽。但这比起\OriginalTerm{圣言}{Word}所做的，从人身上驱逐疾病、魔鬼和死亡本身，又算什么呢？\newline{}4. \Person{狄俄尼索斯}在他们中间被敬拜，因为他教人醉酒；但那真正的救主、万有的主，却因教导节制而被这些人所讥笑。\newline{}5. 不过这些事且放一边。对于祂\OriginalTerm{天主性}{Godhead}的其他奇迹，他们又怎么说呢？有哪个人死亡时，太阳昏暗，大地震动呢？看哪，直到今日，人们仍在死去，古时的人也死去了。在他们身上，何时发生过这样的奇事呢？\newline{}6. 或者，暂且不说祂藉肉身所行的事，只说祂复活以后的事：有谁的教义曾遍及各处，同一无二，从地的一极到另一极，以致他的敬礼如翅膀般飞越每一片土地呢？\newline{}7. 再者，如果基督如他们所说只是人，而非\OriginalTerm{天主圣言}{God the Word}，那么他们所有的神，为何不能阻止祂的敬礼传入他们所在的同一片土地呢？\newline{}8. 相反，为什么\OriginalTerm{圣言}{Word}亲临此地，以祂的教导制止了对他们的敬拜，并使他们的欺诓蒙羞呢？\par

% structure: p0106 source=h2 target=section reason=relative-heading-rank; base=h2

\section{\Emph{50. \OriginalTerm{智者}{Sophists}们的无能和对立因基督的死亡而蒙羞。祂的复活甚至在希腊传说中亦无可比拟。}}

在这人以前，世上已有多位君王和暴君，据记载也有许多智者、术士，出自加色丁人、埃及人和印度人中；试问，他们当中可有谁——不说死后，只论生前——曾有能力如此得胜，竟能用其教导充满全地，并使如此众多的群众脱离偶像的迷信，如同我们的救主引领他们弃绝偶像归向祂自己一样？ 2. 希腊的哲学家们写了许多著作，有似真实之言与辞令技巧；然而，他们展示了什么成果，能如基督的十字架那样伟大呢？他们所教导的精妙理论，至他们死去尚算合理；但即便他们生前似乎拥有的影响力，也受制于彼此间的相互竞争；他们争强好胜，互相攻讦。 3. 但是，天主的圣言——这是最奇妙的事——却以更平凡的语言施教，使那些精选的智者派黯然失色；祂吸引万民归向自己，使他们的学派归于虚无，同时充满了自己的教会；奇妙的是，祂作为人降下、进入死亡，就使那些智者关于偶像的高谈阔论归于虚无。 4. 谁的死亡曾驱逐魔鬼？魔鬼又曾畏惧谁的死亡，如同畏惧基督的死亡？因为哪里呼唤救主的名，那里的所有魔鬼就被驱逐。又有谁曾这样解除人的本性私欲，使淫乱者成为贞洁，使杀人犯不再持剑，使一向被怯懦支配的人成为勇士？ 5. 简言之，除了基督的信德和十字圣号，谁曾说服蛮族之人和各处的异教徒，使他们放下狂妄、追求和平？又有谁曾给予人如此确据的不死不灭，如同基督的十字架和祂身体的复活所给予的？ 6. 希腊人虽讲了各种虚假故事，却不能虚构其偶像的复活——因为他们心中从未想到，死后身体是否可能再度存在。在此处，人们尤其应当接受他们的见证，因为通过这种看法，他们暴露了自己偶像崇拜的弱点，同时为基督留下了可能性，好使祂因此也可在众人中被认识为天主子。\par

% structure: p0108 source=h2 target=section reason=relative-heading-rank; base=h2

\section{五十一、\Emph{节制的新德性。社会革新，由基督宗教净化并带来和平}}

谁是那曾教导童贞，并说明此德在人中并非不可能的？然而基督，我们的救主和万王之王，祂有关此事的教导具有如此能力，连尚未达到法定年龄的儿童，也誓许那超越法律的童贞。 2. 有谁曾能远至斯基泰人、埃塞俄比亚人、波斯人、亚美尼亚人、哥特人，或者那些我们听说在大洋之外或希尔卡尼亚之外的人，甚至埃及人和加色丁人——这些钻研魔法、迷信超乎寻常、习性野蛮的人——中间，去宣讲德行与自制，反对偶像崇拜，如同万有之主，天主的大能，我们的主耶稣基督所做的那样？ 3. 祂不仅藉着自己的门徒宣讲，还将说服带入人心，使人放下残暴的习性，不再事奉祖传的神明，转而认识祂，并藉着祂朝拜圣父。 4. 从前，在偶像崇拜中，希腊人与蛮族彼此征战，对自己的亲属也残忍对待。那时，没有佩剑在手，任何人都无法穿越海洋或陆地，因为他们中间有着永无休止的争斗。 5. 他们整个生活都依靠武器，刀剑对他们如同手杖，是危急时的倚靠；而且，如我前所述，他们仍在事奉偶像，向魔鬼祭献，尽管有这一切偶像崇拜的迷信，他们仍无法被挽回离弃这种风气。 6. 但当他们进入基督的学校之后，奇妙的是，他们如同良心真的被刺痛，便放下了杀人的残暴，不再思念战事：一切都处于和平之中，从今以后，他们所喜爱的，是那些促进友谊的事。\par

% structure: p0110 source=h2 target=section reason=relative-heading-rank; base=h2

\section{五十二、\Emph{魔鬼引发的战争等，由基督宗教平息}}

那么，成就此事的是谁？使彼此相恨之人和睦共处的又是谁？岂不就是父的爱子，众人的共同救主，耶稣基督？祂因自己的爱，为我们的救恩承受了一切。因为自古以来，关于祂要带来的和平，就有预言说，圣经言：\BibleQuote{他们要将刀剑打铸成犁头，把枪矛打铸成镰刀；邦国不再向邦国举剑，也不再学习战事。}\BibleRef{依 2:4} 2. 这至少并非难以置信，因为即使是现在，那些天生习性野蛮的蛮族，只要他们仍向自己国家的偶像祭献，便彼此疯狂，连一刻没有武器都无法忍受： 3. 但当他们听了基督的教导，就立即不再争战，转而从事耕种，不再用武器武装双手，而是举起手来祈祷；一言蔽之，他们不再彼此争斗，从今以后，他们武装自己对抗魔鬼和邪魔，以自制和灵魂的德性战胜他们。 4. 这既是救主神性的一个证明——人从偶像那里不能学到的，却从祂学到了；也是魔鬼和偶像的无能与虚无的明显揭露。因为魔鬼知道自己的无能，为此，他们从前驱使人彼此争战，以免人若停止内讧，便会转而攻击他们。 5. 确实，那些成为基督门徒的人，不再彼此争战，而是以他们的习惯和德行，列阵对抗魔鬼；他们击溃魔鬼，嘲讽他们的首领——魔鬼；以至于青年人有自制，在诱惑中坚忍，在劳苦中坚持，受辱时忍耐，被抢夺时轻看；并且，奇妙的是，他们甚至轻视死亡，成为基督的殉道者。\par

% structure: p0112 source=h2 target=section reason=relative-heading-rank; base=h2

\section{五十三、\Emph{基督秘密地对人的良心说话，一举摧毁整个外邦异教的体系。}}

若论救主的天主性，有一项证据确实极其惊人——哪一个凡人、术士、暴君或君王，曾凭己力与如此众多者交锋，对抗所有偶像崇拜、整个魔鬼军旅、一切法术及希腊人的全部智慧，当他们如此强大、仍处鼎盛且迷惑众人时，一举而遏制他们全体，如同我们的主、真天主圣言所做的那样？祂无形中揭露每个人的错谬，亲身将所有人从这一切中拯救出来，以致崇拜偶像的人现今践踏偶像，以法术闻名的焚烧其书，智慧人则喜爱研究福音胜过一切学问？\newline{}2. 他们先前所崇拜的，现在离弃了；他们先前所讥笑被钉死的，现在却崇拜祂为基督，承认祂是天主。他们中间称为神的，被十字圣号击败，而被钉死的救主则在全世界被宣认为天主、天主子。希腊人崇拜的神明，在他们手中声名扫地，被视为可耻之物；而领受基督教导的人度着比他们更贞洁的生活。\newline{}3. 那么，如果这些事及类似的是人的作为，就让愿意的人指出古人曾有类似的作为，好说服我们。但如果它们证实为，并且的确是，天主的作为而非人的作为，不信者为何如此不虔敬，竟不承认成就这些事的主宰？\newline{}4. 因为他们的情形，正如同一个人看见了受造界的工程，却不认识创造它们的天主。如果他们从祂对宇宙的大能上认识祂的天主性，他们就会知道基督的身体的工程也不是人的，而是万有救主、天主圣言的工程。他们若这样认识，就如保禄所说，\BibleQuote{他们就不会把荣耀的主钉在十字架上了。}\BibleRef{格前 2:8}\par

% structure: p0114 source=h2 target=section reason=relative-heading-rank; base=h2

\section{五十四、\Emph{降生成人的圣言，如同不可见的天主一样，借着祂的工程而为我们所认识。我们借着这些工程认出祂的\OriginalTerm{使人天主化}{deifying}使命。我们满足于列举少数几项，而将众多的辉煌留给愿意瞻望的人。}}

因此，正如一个人若想看见天主——祂本性不可见，绝非肉眼能见——却能从祂的工程中认识和领会祂；同样，理解力不足而不能看见基督的人，至少要从祂身体的工程去领会祂，并试验这些工程是人的还是天主的。\newline{}2. 如果是人的，就让他嘲笑吧；但如果不是人的，而是天主的，就让他承认，并不要嘲笑那不该嘲笑的事；反倒要惊奇，借着如此平常的方式，天主的事竟显示给了我们；借着死亡，不朽竟广施于众人；借着圣言成为人，普遍的天主眷顾及其赐予者和创造者——即天主圣言本身——竟为人所知。\newline{}3. 因为祂成为人，为使我们\OriginalTerm{成为天主}{made God}；祂借着一个身体显示自己，为使我们领悟不可见的父；祂忍受人的侮辱，为使我们继承不朽。因为祂本身丝毫不受伤害，既是不能受苦、不可朽坏的真圣言和天主；而祂在自己毫无痛苦之中，维护并保守了那些正受痛苦、并为他们的缘故忍受这一切的人。\newline{}4. 总之，救主的成就，源于祂成为人，如此的性质和数量，若有人想一一列举，就如同人观看汪洋大海想要数算波浪。人无法尽览全部波浪，因为前来的波浪使试图这样做的人眼花撩乱；同样，想要领会基督在肉身中的全部成就，即使逐一计算，也不可能领会全部，因为超出其思想的比自以为领会的更多。\newline{}5. 所以，最好不要企图述说全部，既然对一部分也难以尽述；只再提一项，而将全部留给您去惊叹。因为一切同样奇妙，人无论转向何方，都能看见圣言的天主性，心生极大敬畏。\par

% structure: p0116 source=h2 target=section reason=relative-heading-rank; base=h2

\section{五十五、\Emph{前文总结。异教神谕的停止等：信仰的传播。真正的君王已来临，使一切篡位者噤声。}}

那么，综上所述，你们理应明白，并以此作为我们以上所述之总结，且应极感惊奇：即自从救主来到我们中间，偶像崇拜不仅不再增长，反而日渐减少，渐趋消亡；希腊人的智慧非但不再进步，反而日渐衰微；至于魔鬼，它们再也不能以幻象、预言和巫术欺骗人，即或胆敢尝试，也被\OriginalTerm{十字圣号}{sign of the Cross}所羞辱。2. 一言以蔽之：请看！救主的教义处处传扬，而一切偶像崇拜及所有与基督信德相悖之事，则日渐消减、失势、没落。你们既看到这一切，就当钦崇那\Quote{超越万有者}、大能的救主，即\OriginalTerm{天主圣言}{God the Word}，并谴责那些被祂击败且除灭的。3. 正如太阳一出，黑暗便不再得势，即或有残余，也必被驱散；同样，如今\OriginalTerm{天主圣言}{the Word of God}的\OriginalTerm{神性显现}{divine Appearing}既已来临，偶像的黑暗便不再得势，普世各方都因祂的训导而蒙光照。4. 又如一位君王统治某国，却不露面，只居守在自己的宫中，多有作乱之徒，趁其退隐之机，擅自称王；他们各自僭取君主之容，欺骗无知之人为王，致使人们只闻其名而不见其人，受其迷惑，或许只因无法进入王宫之故；但当真王显现之时，那些僭妄的骗子便因他的临在而败露，人们既见真王，便离弃先前迷惑他们的人。5. 同样道理，邪神从前惯用的伎俩是欺骗人们，将属天主的尊荣归于自己；但当\OriginalTerm{天主圣言}{the Word of God}以肉体显现，并将祂自己的圣父启示给我们之后，邪神的诡计就终被揭穿并终止，人们转而仰望真天主，即圣父的圣言，从而离弃偶像，得以认识真天主。6. 这正是基督是天主圣言及天主大能的明证。因为人事终必逝去，而基督的圣言却永远长存，众人皆可清楚看出：逝去的是暂时的，而永存者即是天主、即真天主之子、祂的独生圣言。\par

% structure: p0118 source=h2 target=section reason=relative-heading-rank; base=h2

\section{五十六、\Emph{那么，你们要查考圣经，若你们能够，并以此充实此纲要。学习仰望第二次来临和审判。}}

那么，热爱基督的人，愿这就是我们对你所作的奉献，只是一个初步的草图和纲略，简短地叙述基督的信仰及其向我们的神性显现。但你若以此为契机，若你研读圣经的文本，诚心致力其中，你将从圣经更完全、更清晰地学到我们所说的精确细节。2. 因为它们是借着传述天主之事的人，由天主所说所写的。我们则从那些熟悉圣经、并已为基督的神性成为殉道者的受默感教师那里学来的，再转而传授给你这求知的热忱。3. 你还将学到关于祂向我们的第二次光荣而真神性的显现，那时祂不再以卑微，而是以祂自己的光荣——不再以谦恭的容貌，而是以祂自己的威仪——将要来临，不再受苦难，而是从此向众人施予祂自己十字架的果实，即复活和不朽；且不再受审判，而是要审判众人，按各人在肉身内所行的，或善或恶；为善人所存留的是天国，而为行恶者则是永火和黑暗。4. 因为主自己也这样说：\BibleQuote{从此以后你们要看见人子坐在那大能者的右边，乘着天上的云彩降来，在圣父的光荣中。}\BibleRef{玛 26:64} 5. 也正为此，救主也有一句话预备我们迎接那一日，说：\Quote{你们要醒寤，因为他来到的时辰，你们不知道。}因为如真福保禄所说：\Quote{我们众人都必须站在基督的审判台前，为叫各人按在肉身内所行的，或善或恶，领取所应得的。}\par

% structure: p0120 source=h2 target=section reason=relative-heading-rank; base=h2

\section{五十七、\Emph{最重要的是，你们要如此生活，好使你们有权吃这知识树暨生命树，并因此获得永恒的喜乐。颂谢词。}}

然而，为了查考圣经并真实认识它们，需要一种尊荣的生活、纯洁的灵魂，以及那依照基督的德行；如此，理智由这德行引导其道路，就能达成它所渴望的，并在人性能力所能及的范围内，领悟有关天主圣言的事理。2. 因为没有纯洁的心思，没有效法圣人们的生活方式，人就不可能理解圣人们的言语。3. 正如人若愿见太阳的光，必定会擦拭并洁淨自己的眼睛，以某种方式洁净自己，以便能看见所欲见者，好使眼睛变得明亮，从而得见太阳之光；或如人若愿见一城或一地，他必定亲身去到那地观看——同样，那愿领悟谈论天主之事的人，也必须先以其生活方式洗洁和淨化自己的灵魂，并藉着效法圣人们的作为而亲近圣人们自己；这样，在他们共同生活的操练中与他们结合，他也就能明白天主所启示给他们的事，然后，由于与他们紧密相连，就能在审判之日逃避罪人的危险及其火，并领受那为圣人们存留在天国的，\BibleQuote{眼所未见，耳所未闻，人心所未想到的，}\BibleRef{格前 2:9}就是为那些度圣善生活、爱天主圣父、并在我们的主耶稣基督内的人所预备的一切。藉着祂，并同祂，愿尊荣、权能和光荣归于圣父自己，偕同圣子自己，在圣神内，至于无穷之世。阿们。\par

\SourceRecord{Translated by Archibald Robertson.；Nicene and Post-Nicene Fathers, Second Series；Vol. 4.；Edited by Philip Schaff and Henry Wace.；Buffalo, NY: Christian Literature Publishing Co.,；1892.；<http://www.newadvent.org/fathers/2802.htm>.}
